% AxiomForge manuscript — template version 1.0.0.
% Build: latexmk -pdf main.tex   (see README.md)
\documentclass[11pt]{amsart}
\makeatletter
% AxiomForge manuscript preamble — template version 1.0.0.
% This file is copied into every contribution and is owned by it afterwards.
% It fixes the typographic identity of the series; it defines no mathematical notation
% (each paper keeps its own notation in macros.tex).

% ---------------------------------------------------------------- page and fonts
\usepackage[a4paper,hmargin=31mm,top=34mm,bottom=36mm,headsep=10mm,footskip=12mm]{geometry}
\usepackage[T1]{fontenc}
\usepackage[utf8]{inputenc}
\usepackage{lmodern}
\usepackage{amsmath,amssymb,amsthm,mathtools}
\usepackage[final,protrusion=true,expansion=true]{microtype}
\usepackage{booktabs}
\usepackage{array}
\usepackage[shortlabels]{enumitem}
\setlist{itemsep=0.2em,topsep=0.4em}
\usepackage{xcolor}
\linespread{1.04}
\setlength{\emergencystretch}{1.5em}
\allowdisplaybreaks[1]
\numberwithin{equation}{section}
\setcounter{tocdepth}{2}
% Table of contents: sections (and their page numbers) in bold with a little space above
% them, subsections indented in normal weight; after \appendix only the (bold) section
% titles are listed, so the appendices never crowd the contents.
\def\l@section{\@tocline{1}{2.5pt plus 1pt minus 2pt}{0pt}{}{\bfseries\let\@tocpagenum\af@tocpagenumbf}}
\def\af@tocpagenumbf#1{\hss{\bfseries #1}}
\def\l@subsection{\@tocline{2}{0pt}{1.5em}{2.3em}{}}
\g@addto@macro\appendix{\addtocontents{toc}{\protect\setcounter{tocdepth}{1}}}
% Every appendix starts on a new page; the unnumbered back matter (\afbacksection) keeps the
% original \section and flows normally after the last appendix.
% (Local implementation through a hook, so that \section itself is not redefined.)
\newif\ifaf@appendixpage
\AddToHook{cmd/section/before}{\ifaf@appendixpage\clearpage\fi}
\g@addto@macro\appendix{\global\af@appendixpagetrue}

% ---------------------------------------------------------------- links
\definecolor{afblue}{rgb}{0.08,0.18,0.42}
\usepackage[pdfusetitle,bookmarks=true,bookmarksnumbered=true,bookmarksopen=true,
  colorlinks=true,linkcolor=afblue,citecolor=afblue,urlcolor=afblue,
  pdfpagelabels=true,pdfdisplaydoctitle=true]{hyperref}
\usepackage[capitalise,noabbrev]{cleveref}
% Local addition: unique PDF anchors for theorem-like environments numbered within sections.
\hypersetup{hypertexnames=false}

% ---------------------------------------------------------------- theorem environments
% thmtools keeps cleveref names correct for environments sharing one counter.
\usepackage{thmtools}
\declaretheorem[style=plain,numberwithin=section,name=Theorem]{theorem}
\declaretheorem[style=plain,sibling=theorem,name=Proposition]{proposition}
\declaretheorem[style=plain,sibling=theorem,name=Lemma]{lemma}
\declaretheorem[style=plain,sibling=theorem,name=Corollary]{corollary}
\declaretheorem[style=plain,numbered=no,name=Theorem]{theorem*}
\declaretheorem[style=definition,sibling=theorem,name=Definition]{definition}
\declaretheorem[style=definition,sibling=theorem,name=Example]{example}
\declaretheorem[style=remark,sibling=theorem,name=Remark]{remark}
\crefname{equation}{}{}
\Crefname{equation}{Equation}{Equations}

% ---------------------------------------------------------------- series metadata
% \afidentifier{AF-0001}, \afversion{1.0.0}, \afdate{25 November 2026},
% \afrepository{https://github.com/axiomforge-math/mathematics/tree/main/results/...},
% \afdisclosure{short AI-assisted research disclosure},
% \afstatus{verification status} (local addition; optional)
\newcommand{\afidentifier}[1]{\gdef\af@identifier{#1}}
\newcommand{\afversion}[1]{\gdef\af@version{#1}}
\newcommand{\afdate}[1]{\gdef\af@date{#1}}
\newcommand{\afrepository}[1]{\gdef\af@repository{#1}}
\newcommand{\afdisclosure}[1]{\gdef\af@disclosure{#1}}
\newcommand{\afcorrespondence}[1]{\gdef\af@correspondence{#1}}
\gdef\af@correspondence{}
% Local addition: verification status of the contribution, printed as a 'Status' row.
\newcommand{\afstatus}[1]{\gdef\af@status{#1}}
\gdef\af@status{}

% The publication block printed under the abstract on the title page.
\newcommand{\afpublicationblock}{%
  \begingroup\footnotesize\noindent
  \begin{tabular}{@{}l@{\quad}p{0.72\textwidth}@{}}
    \textit{Identifier} & \af@identifier\quad(version \af@version, \af@date)\\[1pt]
    \ifx\af@status\@empty\else\textit{Status} & \af@status\\[1pt]\fi
    \textit{Repository} & \url{\af@repository}\\[1pt]
    \ifx\af@correspondence\@empty\else\textit{Correspondence} & \af@correspondence\\[1pt]\fi
    \textit{Disclosure} & \af@disclosure
  \end{tabular}\par
  \endgroup}

% Front matter: title page, then contents, then the body on a fresh page.
\newcommand{\affrontmatter}{%
  \maketitle
  \vspace*{\fill}
  \afpublicationblock
  \thispagestyle{plain}%
  \clearpage
  {\hypersetup{linkcolor=black}\fontsize{9}{10.6}\selectfont\tableofcontents}
  \clearpage}

% Unnumbered back-matter sections (amsart lists starred sections in the contents).
\newcommand{\afbacksection}[1]{\global\af@appendixpagefalse\section*{#1}}

\makeatother
% Notation local to this paper.
\newcommand{\N}{\mathbb{N}}
\newcommand{\Z}{\mathbb{Z}}
\newcommand{\Q}{\mathbb{Q}}
\newcommand{\R}{\mathbb{R}}

% special functions and the target numbers
\DeclareMathOperator{\Ti}{Ti}
\newcommand{\Li}{\operatorname{Li}}

% general operators
\DeclareMathOperator*{\Res}{Res}
\DeclareMathOperator{\diag}{diag}
\newcommand{\pos}[1]{\left(#1\right)_{+}}
\newcommand{\floor}[1]{\left\lfloor #1\right\rfloor}
\newcommand{\vG}[1]{v_{#1}^{\mathrm{G}}}
\newcommand{\abs}[1]{\left\lvert #1\right\rvert}
\newcommand{\one}{\mathbf{1}}

% objects of the construction
\newcommand{\Rsp}{\mathcal{R}_K}              % rational functions with poles at the x_k
\newcommand{\mom}{\mathcal{S}_w}              % rational moment sum
\newcommand{\nrm}{\mathsf{S}}                 % normalising factor
\newcommand{\Pa}{\mathcal{A}}                 % parity sum A(t)
\newcommand{\rest}[1]{#1\langle z\rangle}     % restricted power series
\newcommand{\aconst}{\mathfrak{a}}            % arithmetic constant
\newcommand{\bconst}{\mathfrak{b}}            % analytic constant
\newcommand{\Iprof}{I}                        % ordinary-prime profile integral
\newcommand{\Dk}{D_\kappa}                    % (2-kappa)^2/4
\newcommand{\bvec}{\mathbf{b}}                % binomial column vector
\newcommand{\vvec}{\mathbf{v}}                % monomial column vector
\newcommand{\diff}{\,\mathrm{d}}

% title pieces (robust, so that amsart's upper-casing of the title leaves them unchanged)
\DeclareRobustCommand{\titlechi}{$\chi_2(a/q)$}
\DeclareRobustCommand{\titleti}{$\mathrm{Ti}_2(a/q)$}
\DeclareRobustCommand{\titlea}{$a=3,4,5,8$}


\title[Irrationality of chi and Ti at the numerators 3, 4, 5, 8]%
{\texorpdfstring{Irrationality of \titlechi{} and \titleti{} for the numerators \titlea}%
{Irrationality of chi2(a/q) and Ti2(a/q) for the numerators a = 3, 4, 5, 8}}
\hypersetup{pdftitle={Irrationality of χ₂(a/q) and Ti₂(a/q) for the numerators a = 3, 4, 5, 8}}
\author{AxiomForge}
\urladdr{https://github.com/axiomforge-math}
\subjclass[2020]{Primary 11J72; Secondary 33B30}
\keywords{Legendre chi function, inverse tangent integral, irrationality, Hankel determinants,
Pad\'e-type approximation, $p$-adic valuations, logarithmic energy}
\date{26 November 2026}

\afidentifier{AF-0004}
\afversion{0.1.0}
\afdate{26 November 2026}
\afrepository{https://github.com/axiomforge-math/mathematics/tree/main/results/AF-0004-legendre-chi-general-numerator-irrationality}
\afcorrespondence{\href{mailto:axiomforge.math@gmail.com}{axiomforge.math@gmail.com}}
\afstatus{Preprint --- complete written argument, not yet independently verified; not reviewed by an independent human expert. The finite computations are reproduced by separately written checkers (AI-internal). Lean formalization in preparation: the main theorems are proved in Lean only conditionally on the estimates.}
\afdisclosure{AI-assisted research. The candidate arguments and the original certificate programs were produced with GPT-6.2 XHigh (web version); their reconstruction, the check computations, this manuscript and the Lean formalization with Claude Sonnet~5.5, Claude Fable~5.2 and Claude Opus~5.5 (Anthropic), under the direction of AxiomForge. No independent human review has taken place. See the AI-assisted research disclosure at the end.}

\begin{document}

\begin{abstract}
Let $\chi_2(x)=\sum_{j\ge0}x^{2j+1}/(2j+1)^2$ be Legendre's chi function and
$\Ti_2(x)=\sum_{j\ge0}(-1)^jx^{2j+1}/(2j+1)^2$ the inverse tangent integral. We present a
candidate proof that, for odd $q$ prime to $a$, the number $\chi_2(a/q)$ is irrational when
$(a,q)$ satisfies $a=3$, $q\ge131$; $a=4$, $q\ge995$; $a=5$, $q\ge829$; or $a=8$, $q\ge12061$;
and that $\Ti_2(a/q)$ is irrational when $a=3$, $q\ge257$; $a=4$, $q\ge995$; $a=5$,
$q\ge1609$; or $a=8$, $q\ge12061$. The same holds at $-a/q$. These ranges begin well below
the range $q>54e^4a^2$ that we read off from a general criterion for shifted polylogarithms. The
rational approximations are Hankel determinants of a moment functional with poles at the
negative half-integers, with more poles than rows, in the ratio $7/6$ or $5/4$; this is the
construction of an earlier AxiomForge preprint on the numerator $2$, which itself starts from
work of Krisman. The numerators treated here need three arithmetic estimates that are absent
for the numerator $2$: a bound at the primes dividing $a$, where each pole can be used only
once; a bound at the prime $2$ when $a$ is odd, by forward differences in a basis adapted to
the poles; and the parity-minor bound at $2$ for $a=4$ and $a=8$. A single inequality carries
each case from its threshold to all larger denominators. The finite inequalities on which the
argument rests have been recomputed by separately written programs. The argument has not been
checked by an independent human expert and its main theorems are not formally verified, so the
result should be read as a candidate.
\end{abstract}

\affrontmatter

\section{Introduction}\label{sec:introduction}

For $\abs{x}\le 1$ let
\[
  \chi_2(x)=\sum_{j\ge0}\frac{x^{2j+1}}{(2j+1)^2}=\tfrac12\bigl(\Li_2(x)-\Li_2(-x)\bigr),
  \qquad
  \Ti_2(x)=\sum_{j\ge0}\frac{(-1)^jx^{2j+1}}{(2j+1)^2}=\operatorname{Im}\Li_2(ix)
\]
be Legendre's chi function and the inverse tangent integral. This paper presents a proof of the
following statement.

\begin{theorem}\label{thm:main}
Let $a\in\{3,4,5,8\}$ and let $q$ be an odd integer prime to $a$.
\begin{enumerate}[(i)]
\item The number $\chi_2(a/q)$ is irrational if $q\ge q_+(a)$, where
$q_+(3)=131$, $q_+(4)=995$, $q_+(5)=829$ and $q_+(8)=12061$.
\item The number $\Ti_2(a/q)$ is irrational if $q\ge q_-(a)$, where
$q_-(3)=257$, $q_-(4)=995$, $q_-(5)=1609$ and $q_-(8)=12061$.
\end{enumerate}
Since $\chi_2$ and $\Ti_2$ are odd, the same holds for $\chi_2(-a/q)$ and $\Ti_2(-a/q)$.
\end{theorem}

The thresholds are inclusive, and they are sufficient for our estimates; we do not claim that
they are optimal (\cref{sec:limitations}).

\emph{Status.} The proof is a candidate. It was reconstructed and checked by AI systems from
candidate arguments produced by another AI system (see the disclosure at the end); no
independent human expert has reviewed it. In Lean only the final deduction is formalised, and
only conditionally on the three estimates of \cref{thm:three-estimates}. All finite
inequalities used in the proof have been recomputed by programs written separately from those
that first produced them (Appendices~\ref{app:certificates} and~\ref{app:consistency}); this
checks the numbers, not the asymptotic arguments. \Cref{subsec:step-status} gives the status of
each step, and \cref{rem:scrutiny} lists the points that most need independent scrutiny.

\subsection{Background}\label{subsec:background}

If $q$ is even or not prime to $a$, the fraction $a/q$ reduces to one with a smaller numerator
or an even numerator and odd denominator; we treat the four numerators $3,4,5,8$ with odd
denominators prime to them. The irrationality of $\chi_2$ and $\Ti_2$ at rational points with
large denominator follows from linear independence criteria for shifted polylogarithms. Hata's
measures~\cite{Hata1990} are stated at reciprocals of integers. The criterion of David,
Hirata-Kohno and Kawashima~\cite{DHK2022Acta} applies to the shifted dilogarithm
$\sum_{n\ge1}x^n/(n-\frac12)^2$ at rational points; Krisman evaluates it at $x=4/q^2$ and
obtains the irrationality of $\chi_2(2/q)$ for $q\ge11\,795$ \cite[\S1.2]{Krisman2026chi2}.
The same evaluation at $x=\pm a^2/q^2$ gives the range $q>54e^4a^2$, that is, $q$ larger than
about $2.7\cdot10^4$, $4.7\cdot10^4$, $7.4\cdot10^4$ and $1.9\cdot10^5$ for $a=3,4,5,8$. We have not re-derived this range from
\cite{DHK2022Acta}; we quote it only to indicate the order of magnitude of what was known, and
we have not found explicit smaller thresholds for these numerators in the literature.

For small denominators the closest works are Krisman's preprint~\cite{Krisman2026chi2}, which
proves that $\chi_2(1/q)\notin\Q$ for $q\ge3$, $q\ne4$, that $\Ti_2(1/q)\notin\Q$ for $q\ge5$,
and that $\chi_2(2/q)\notin\Q$ for odd $q\ge175$, and the AxiomForge
preprint~\cite{AxiomForge2026numeratorTwo}, a candidate proof that $\chi_2(2/q)$ and
$\Ti_2(2/q)$ are irrational for odd $q\ge81$. Both are unrefereed and AI-assisted. We use none
of their theorems as an input: every estimate needed here is proved in this paper, including
those that are restatements, for general numerator, of estimates
of~\cite{AxiomForge2026numeratorTwo}. The method is a Hankel-determinant construction in the
spirit of Zudilin's determinantal approach~\cite{Zudilin2017}: Krisman takes the Hankel
determinants of the functional $g\mapsto\sum_{t\ge1}w^t\,(g(u)\,u/D(u))'|_{u=t}$, where
$D(u)=\prod_{k<h}(u+k+\frac12)$ has as many half-integer poles as the size $h$ of the
determinant; in~\cite{AxiomForge2026numeratorTwo} the number of poles exceeds $h$ in the ratio
$7/6$.

\subsection{The construction}\label{subsec:constructions}

Fix an integer $a\ge1$, an odd integer $q>a$ prime to $a$ and a sign $\sigma\in\{1,-1\}$, and
put
\begin{equation}\label{eq:basic-notation}
  y=\frac qa,\qquad w=\frac{\sigma}{y^2}=\frac{\sigma a^2}{q^2},\qquad W=\abs w,\qquad
  \xi=\sum_{j\ge0}\frac{\sigma^jy^{-2j-1}}{(2j+1)^2},
\end{equation}
so that $\xi=\chi_2(a/q)$ for $\sigma=1$ and $\xi=\Ti_2(a/q)$ for $\sigma=-1$. For integers
$K\ge h\ge1$ let
\[
  x_k=-k-\tfrac12\ (0\le k<K),\qquad D_K(u)=\prod_{k<K}(u-x_k),\qquad
  \varphi(u)=\frac{u}{D_K(u)}.
\]
In \cref{sec:pencils} we define an affine functional $\mu_X$ on the rational functions with
simple poles at the $x_k$ such that $\mu_\xi(R)=\sum_{m\ge1}w^mR'(m)$, and the polynomial
\[
  T(X)=\nrm\cdot\det\bigl[\mu_X(\varphi(u)\,u^{i+j})\bigr]_{0\le i,j<h}\in\Q[X],
  \qquad \nrm=\frac{K!^{\,h}}{\prod_{i<h}i!^2},
\]
of degree at most $h$. We use two families of sizes, indexed by $n\ge1$:
\[
  (h,K)=(6n,7n)\quad(\kappa=\tfrac76),\qquad (h,K)=(4n,5n)\quad(\kappa=\tfrac54),
\]
where $\kappa=K/h$. The first is used for $\chi_2(3/q)$ and the second for the seven other
cases of \cref{thm:main}. Throughout,
\[
  \Dk=\frac{(2-\kappa)^2}{4},\qquad\text{so}\qquad D_{7/6}=\frac{25}{144},\quad
  D_{5/4}=\frac9{64}.
\]

\subsection{The argument in one page}\label{subsec:one-page}

For a prime $p$ and a nonzero $F\in\Q[X]$ let $\vG p(F)$ be the minimum of the $p$-adic
valuations of the coefficients of $F$. Let $\ell=\prod_pp^{-\vG p(T)}$, so that $\ell T\in\Z[X]$.
If $\xi=r/s$ were rational, then $s^h\ell T(r/s)$ would be an integer. The proof consists in
showing that this integer is nonzero for infinitely many $n$, while its absolute value is at
most $\exp\bigl(h^2(\bconst-\aconst+o(1))\bigr)$ with $\bconst<\aconst$
(\cref{prop:criterion}). This needs three estimates.

\begin{theorem}\label{thm:three-estimates}
Let $a\ge1$, let $q>a$ be odd and prime to $a$, let $\sigma=\pm1$ and $\kappa\in\{7/6,5/4\}$,
and let $T=T_n$ be the polynomial attached to the $n$-th size of ratio $\kappa$, with $h=h_n$.
Write $a=2^ea_0$ with $a_0$ odd.
\begin{enumerate}[(i)]
\item \emph{Arithmetic.} Along the $n$ with $T_n\ne0$,
\begin{multline*}
  \liminf\frac{1}{h^2}\sum_p\vG p(T)\log p\ \ge\ \aconst_\sigma(a,q)\\
  :=\Iprof_\kappa+c_2\log2-(2\kappa-1)\log a_0+\frac13\log q-\Dk\log(q^2-\sigma a^2),
\end{multline*}
where $\Iprof_{7/6}=-2.113164137$, $\Iprof_{5/4}=-2.355639833$, and
\[
  c_2=\begin{cases}2(\kappa-1)+2\Dk&\text{if $a$ is odd and }\sigma=1,\\
  \kappa-1&\text{if $a$ is odd and }\sigma=-1,\\
  C_e(\kappa)=\min_{0\le x\le1}f_e(x)&\text{if }e\ge1,\end{cases}
\]
with $f_e(x)=2(\kappa-1)-2(e+1)\kappa x+(e+2)x^2+(\frac e2+2)(1-x)^2$.
\item \emph{Size.} For every comparison density $\rho$ and every $\lambda$ satisfying the tail
condition of \cref{lem:tail}, $\limsup h^{-2}\log\abs{T(\xi)}\le\bconst(\rho,\lambda)$, an
explicit number that depends on $(a,q)$ only through $W$.
\item \emph{Nonvanishing.} Let $r\in\Q$. If $q\ge14a$, then $T_n(r)\ne0$ for every
sufficiently large prime $n$.
\end{enumerate}
\end{theorem}

\begin{proof}
Part~(i) is \cref{thm:arithmetic}, part~(ii) is \cref{thm:decay}, and part~(iii) is
\cref{thm:nonvanishing} together with \cref{prop:local-dets}, since $q\ge14a$ gives
$\abs{\sigma q^2/a^2}\ge196>189$.
\end{proof}

\Cref{tab:constants} lists the constants at the six thresholds. For $a=4$ and $a=8$ the two
signs are treated together, with the smaller arithmetic constant $\aconst_-$
(\cref{rem:worse-sign}).

\begin{table}[ht]
\centering
\caption{The six threshold comparisons, per $h^2$. The arithmetic constant is
$\aconst_\sigma(a,q_0)$ of \cref{thm:three-estimates}, with the value of $c_2$ shown; $\bconst$ is
the certified size exponent (Appendix~\ref{app:decay-certificates}). Decimals are rounded in the
safe direction.}
\label{tab:constants}
\small
\begin{tabular}{@{}llllllll@{}}
\toprule
$a$ & $q_0$ & function & $\kappa$ & $c_2$ & $\aconst\ge$ & $\bconst\le$ & $\aconst-\bconst\ge$\\
\midrule
$3$ & $131$ & $\chi_2$ & $7/6$ & $49/72$ & $-3.173875358$ & $-3.187821144$ & $0.013945786$\\
$3$ & $257$ & $\Ti_2$ & $5/4$ & $1/4$ & $-3.541276253$ & $-3.546389753$ & $0.005113500$\\
$4$ & $995$ & both & $5/4$ & $-337/112$ & $-4.081754611$ & $-4.083454335$ & $0.001699724$\\
$5$ & $829$ & $\chi_2$ & $5/4$ & $25/32$ & $-3.878258885$ & $-3.879559185$ & $0.001300300$\\
$5$ & $1609$ & $\Ti_2$ & $5/4$ & $1/4$ & $-4.211960841$ & $-4.212674322$ & $0.000713481$\\
$8$ & $12061$ & both & $5/4$ & $-9/2$ & $-4.985336979$ & $-4.986095509$ & $0.000758530$\\
\bottomrule
\end{tabular}
\end{table}

The estimates are obtained as follows.
\begin{itemize}
\item \emph{Prime $2$, $a$ odd} (\cref{subsec:two-odd}). Here $y$ is a $2$-adic unit, and $2$
divides both the denominators of the poles and $q^2-\sigma a^2$. In bases of polynomials that
vanish at the first and at the last poles, each entry of the matrix is a forward difference of a
sequence obeying a first-order recursion, and its valuation is bounded directly.
\item \emph{Prime $2$, $a$ even} (\cref{subsec:two-even}). In the coordinate $z=2u$ the poles
are odd and the summation points even, and a Vandermonde-type minor whose nodes share a parity
is divisible by about $2^{t^2}$; the Cauchy--Binet formula reduces $\vG2(T)$ to the minimisation
of $f_e$ in the proportion $x$ of pole terms.
\item \emph{Odd primes dividing $a$} (\cref{subsec:p-divides-a}). The pole $x_k$ carries the
factor $y^{2k+1}$, of valuation $-(2k+1)v_p(a)$, but it enters the matrix through a term of rank
one, so a term of the determinant uses each pole at most once; the $h$ most expensive poles cost
$(2\kappa-1)v_p(a)h^2$.
\item \emph{Odd primes dividing $q(q^2-\sigma a^2)$} (\cref{sec:exceptional}). For $p\mid q$,
reflecting $m\mapsto-m$ makes the series $p$-adically convergent, and each level offers at most
three terms of valuation at least $2v_p(q)m$, which gives $\frac13v_p(q)h^2$; for
$p\mid q^2-\sigma a^2$ the only loss comes from the denominators $(1-w)^{b+1}$, and a convexity
argument gives $\Dk\,v_p(q^2-\sigma a^2)h^2$.
\item \emph{Ordinary primes} (\cref{sec:ordinary}). For $p\asymp h$ two lower bounds, by Taylor
jets in each residue class modulo $p$ and by global partial fractions with Vandermonde gains,
are small minimisation problems over the residue classes; their maximum, integrated against
$\log p$ by the prime number theorem, gives $\Iprof_\kappa$, and the integral is evaluated
exactly (Appendix~\ref{app:profile-certificate}).
\item \emph{Size} (\cref{sec:decay}). The determinant at $X=\xi$ is a sum over configurations
of $h$ points on the summation sites; Vandermonde bounds and a discrete energy inequality
compare it with an explicit step density, and the resulting number $\bconst$ is evaluated in
ball arithmetic (Appendix~\ref{app:decay-certificates}).
\item \emph{Nonvanishing} (\cref{sec:nonvanishing}). For $h=tp$ with $p$ prime, a basis adapted
to the residue classes modulo $p$ makes the scaled determinant block triangular modulo $p$, with
diagonal blocks equal, up to units, to three fixed local Hankel determinants, which are
nonzero.
\item \emph{Assembly} (\cref{sec:proofs}). The margins of \cref{tab:constants} are positive,
and for $q\ge q_0\ge2a$ the margin grows by at least $\frac{10}{27}\log(q/q_0)$
(\cref{prop:propagation}), which carries each threshold to all larger admissible $q$.
\end{itemize}

\subsection{What is new and what is inherited}\label{subsec:new}

The functional, the criterion, the reflection at primes dividing $q$, the size bound with its
shift to smaller $\abs w$, and the strategy of proving nonvanishing modulo a prime proportional
to $h$ go back to Krisman~\cite{Krisman2026chi2}. The ratio $K/h>1$, the parity-minor bound at
$2$ for $a=2$, the residue-class profile for the ordinary primes and the Taylor-jet form of the
nonvanishing argument are those of~\cite{AxiomForge2026numeratorTwo}, where $\kappa=7/6$. We
reprove all of them for general $a$ and both ratios; in
\cref{sec:ordinary,sec:decay,sec:nonvanishing} the numerator enters only through $w$, and the
proofs are those of~\cite{AxiomForge2026numeratorTwo} with a simpler weight. The following is
new.
\begin{enumerate}[(a)]
\item \emph{Primes dividing the numerator} (\cref{prop:p-divides-a}), which do not occur for
$a\le2$ apart from the prime $2$.
\item \emph{The prime $2$ for odd $a$} (\cref{thm:two-odd}). Its cost depends on
$\nu=v_2(q^2-\sigma a^2)$, but the dependence cancels against the contribution of $2$ to
$\log(q^2-\sigma a^2)$, so that $\aconst_\sigma(a,q)$ contains the fixed multiples
$\frac{49}{72}\log2$ or $\frac{25}{32}\log2$ ($\sigma=1$) and $\frac16\log2$ or $\frac14\log2$
($\sigma=-1$).
\item \emph{The prime $2$ for $a=2^ea_0$ with $e\ge2$} (\cref{thm:two-even}): the parity-minor
argument of~\cite{AxiomForge2026numeratorTwo} with pole costs $e(2k+1)$ and site costs $2em$.
For $e=3$ and $\kappa=5/4$ the minimum of $f_e$ is attained at $x=1$ and equals $-e(2\kappa-1)$,
the value given by the rank-one argument of~(a).
\item \emph{The ratio $5/4$}: a second profile integral, a second triple of local determinants
(of size $4$), and five new comparison densities.
\item \emph{Propagation in $q$ for general $a$} (\cref{prop:propagation}).
\end{enumerate}

\subsection{Organisation and notation}\label{subsec:notation}

\Cref{sec:pencils} defines the functional, the determinants and the criterion.
\Cref{sec:prime-two,sec:exceptional,sec:ordinary} bound the denominators, \cref{sec:decay} the
size, and \cref{sec:nonvanishing} proves nonvanishing; \cref{sec:proofs} assembles the proof and
records the status of each step, and \cref{sec:limitations} discusses what the thresholds do and
do not say. Appendix~\ref{app:certificates} describes the finite computations and
Appendix~\ref{app:consistency} compares the estimates with exact computations at small sizes.
\Cref{tab:notation} lists the notation used throughout.

\begin{table}[ht]
\centering
\caption{Notation.}\label{tab:notation}
\small
\begin{tabular}{@{}p{0.27\textwidth}p{0.66\textwidth}@{}}
\toprule
$a$, $q$, $\sigma$; $y$, $w$, $W$, $\xi$ & numerator, odd denominator prime to $a$, sign $\pm1$;
see~\eqref{eq:basic-notation}\\
$e$, $a_0$; $Q$, $\nu$ & $a=2^ea_0$ with $a_0$ odd; $Q=\sigma y^2=1/w$, $\nu=v_2(Q-1)$ for odd
$a$\\
$\tau_k$ & partial sum $\sum_{j\le k}\sigma^jy^{-2j-1}/(2j+1)^2$ of $\xi$, \eqref{eq:tau}\\
$h$, $K$; $\kappa$, $\Dk$ & size, number of poles; $K/h\in\{\frac76,\frac54\}$,
$(2-\kappa)^2/4$\\
$x_k$, $D_K$, $\varphi$, $\Phi$ & $-k-\frac12$; $\prod_{k<K}(u-x_k)$; $u/D_K$; $K!\,\varphi$\\
$\Omega_b$, $\mom$, $\mu_X$ & $w^{\max(b,1)}(1-w)^{-b-1}$; rational moment sum; the functional
\eqref{eq:functional}\\
$\Delta$, $\nrm$, $T$ & Hankel determinant, normaliser, $T=\nrm\Delta$, see~\eqref{eq:Delta-T}\\
$\beta_i$, $\vvec(u)$, $\bvec(u)$ & $\binom ui$; the columns $(u^i)_{i<h}$ and
$(\beta_i(u))_{i<h}$\\
$\vG p$ & Gauss valuation of a polynomial in $X$ over $\Q_p$\\
$\aconst_\sigma(a,q)$, $\bconst$, $\Iprof_\kappa$ & arithmetic constant, size exponent,
ordinary-prime constant\\
$\psi(t)$, $\pos x$, $s_2(n)$ & $t\log\abs t-t$ with $\psi(0)=0$; $\max(x,0)$; binary digit
sum, so $v_2(n!)=n-s_2(n)$\\
\bottomrule
\end{tabular}
\end{table}

\emph{Conventions.} $p$ is always a prime and $v_p$ the $p$-adic valuation on $\Q_p$. The Gauss
valuation satisfies $\vG p(FG)=\vG p(F)+\vG p(G)$ and $\vG p(F+G)\ge\min(\vG p(F),\vG p(G))$, with
$\vG p(0)=+\infty$, and is invariant under $X\mapsto X+c$ for $c\in\Z_p$; hence the determinant of
a matrix over $\Q_p[X]$ has Gauss valuation at least the minimum over the terms of its Leibniz
expansion. For $u$ in any $\Q$-algebra, $\binom uk=u(u-1)\cdots(u-k+1)/k!$. The implied constants
in $O(\cdot)$ may depend on $a$, $q$ and $\kappa$, but not on $n$.

\section{The half-integer pencils}\label{sec:pencils}

This section defines the objects of the proof: the functional $\mu_X$, the polynomial $T$, its
expansion over the summation points, a rank-one expansion used three times, and the
irrationality criterion. All later sections are
statements about $T$, and the expansion~\eqref{eq:real-sites} with its $p$-adic analogues is the
starting point of \cref{sec:prime-two,sec:exceptional,sec:ordinary,sec:decay}.

Throughout, $a\ge1$ is an integer, $q>a$ is odd and prime to $a$, $\sigma=\pm1$, and
$y,w,W,\xi$ are as in~\eqref{eq:basic-notation}. We put
\begin{equation}\label{eq:tau}
  \tau_k=\sum_{j=0}^{k}\frac{\sigma^jy^{-2j-1}}{(2j+1)^2}\qquad(k\ge0),
\end{equation}
the partial sums of the series defining $\xi$.

\subsection{Moment sums and the functional}\label{subsec:moment-sums}

For $b\ge0$ let $\Omega_b(w)=w^{\max(b,1)}(1-w)^{-b-1}\in\Q(w)$. For a polynomial $g\in\Q[u]$,
expanded in the Newton basis as $g(u)=\sum_b\Delta^bg(0)\binom ub$, put
\[
  \mom(g)=\sum_{b\ge0}\Delta^bg(0)\,\Omega_b(w).
\]

\begin{lemma}\label{lem:moment-sums}
Let $g\in\Q[u]$ and $w\in\Q\setminus\{0,1\}$.
\begin{enumerate}[(a)]
\item The identity $\sum_{m\ge1}g(m)w^m=\mom(g)$ holds in $\Q[[w]]$. Consequently, for every
absolute value $\abs\cdot_v$ on $\Q$ with $\abs w_v<1$, the series $\sum_{m\ge1}g(m)w^m$
converges in the completion $\Q_v$ to $\mom(g)$.
\item The identity $\mom(g)=-\sum_{m\ge0}g(-m)w^{-m}$ holds in $\Q[[w^{-1}]]$, and hence in
$\Q_v$ whenever $\abs w_v>1$.
\item Let $p$ be a prime such that $v_p(w)\ge0$ and $1-w$ is a $p$-adic unit. If
$g(\Z)\subset\Z_p$, then $\mom(g)\in\Z_p$.
\end{enumerate}
\end{lemma}

\begin{proof}
(a) For $b\ge1$, $\sum_{m\ge1}\binom mbw^m=w^b(1-w)^{-b-1}$, and $\sum_{m\ge1}w^m=w(1-w)^{-1}$.
(b) We have $\binom{-m}{b}=(-1)^b\binom{m+b-1}{b}$, so for $b\ge1$, with $v=w^{-1}$,
\[
 \sum_{m\ge0}\binom{-m}bv^m=(-1)^b\frac{v}{(1-v)^{b+1}}=(-1)^b\frac{w^{b}}{(w-1)^{b+1}}
 =-\Omega_b(w),
\]
and $\sum_{m\ge0}v^m=w/(w-1)=-\Omega_0(w)$. (c) The Newton coefficients $\Delta^bg(0)$ are
integral combinations of $g(0),\dots,g(b)$, and $\Omega_b(w)\in\Z_p$ because $w\in\Z_p$ and
$1-w\in\Z_p^\times$.
\end{proof}

Let $K\ge1$, $x_k=-k-\frac12$ for $0\le k<K$, and $D_K(u)=\prod_{k<K}(u-x_k)$. Let $\Rsp$ be
the $\Q$-vector space spanned by $\Q[u]$ and the fractions $1/(u-x_k)$, $0\le k<K$; every
rational function whose only poles are simple poles at some of the $x_k$ lies in $\Rsp$, and
its decomposition into a polynomial part and principal parts is unique. For an indeterminate
$X$ define the $\Q$-linear map $\mu_X\colon\Rsp\to\Q[X]$ by
\begin{equation}\label{eq:functional}
  \mu_X(R)=\mom(R')\quad(R\in\Q[u]),\qquad
  \mu_X\Bigl(\frac1{u-x_k}\Bigr)=-4\sigma^ky^{2k+1}\,(X-\tau_k).
\end{equation}
Each value is a polynomial of degree at most one in $X$.

\begin{lemma}[pole series]\label{lem:pole-series}
For every $R\in\Rsp$ the series $\sum_{m\ge1}w^mR'(m)$ converges absolutely and equals
$\mu_\xi(R)$.
\end{lemma}

\begin{proof}
For polynomials this is \cref{lem:moment-sums}(a) with the real absolute value. For
$R=1/(u-x_k)$ we have $R'(m)=-4/(2m+2k+1)^2$. Writing $2m+2k+1=2j+1$ with $j=m+k\ge k+1$
and $w^m=\sigma^{j-k}y^{-2(j-k)}$,
\[
  \sum_{m\ge1}w^mR'(m)=-4\sigma^ky^{2k+1}\sum_{j>k}\frac{\sigma^jy^{-2j-1}}{(2j+1)^2}
  =-4\sigma^ky^{2k+1}(\xi-\tau_k).\qedhere
\]
\end{proof}

The same computation will be used $2$-adically in \cref{sec:prime-two} and $p$-adically in
\cref{sec:exceptional}. Lemma~\ref{lem:pole-series} is Krisman's pole series
\cite[Lemma~2.1]{Krisman2026chi2}, with an arbitrary number $K$ of poles and an arbitrary
numerator $a$.

\subsection{The determinants}\label{subsec:determinants}

Let $K\ge h\ge1$ be integers, and put
\begin{equation}\label{eq:phi}
  \varphi(u)=\frac{u}{D_K(u)}.
\end{equation}
For every polynomial $g$ the product $\varphi g$ lies in $\Rsp$. Define
\begin{equation}\label{eq:Delta-T}
  \Delta(X)=\det\bigl[\mu_X(\varphi(u)u^{i+j})\bigr]_{0\le i,j<h},\qquad
  \nrm=\frac{K!^{\,h}}{\prod_{i<h}i!^2},\qquad T(X)=\nrm\,\Delta(X).
\end{equation}
These are polynomials in $X$ with rational coefficients and degree at most $h$, and
$\Delta(\xi)$ is the determinant of the real matrix
$\bigl[\sum_{m\ge1}w^m(\varphi u^{i+j})'(m)\bigr]$ by \cref{lem:pole-series}.

\begin{lemma}[change of basis]\label{lem:change-of-basis}
Let $f_0,\dots,f_{h-1}$ and $g_0,\dots,g_{h-1}$ be polynomials of degree less than $h$ with
coefficient matrices $C_f$, $C_g$ in the monomial basis. Then
$\det[\mu_X(\varphi f_ig_j)]=\det C_f\det C_g\,\Delta(X)$. In particular, with
$\beta_i(u)=\binom ui$ and
\begin{equation}\label{eq:Phi}
  \Phi(u)=K!\,\varphi(u)=\frac{K!\,u}{D_K(u)},
\end{equation}
one has $T(X)=\det[\mu_X(\Phi\beta_i\beta_j)]_{0\le i,j<h}$.
\end{lemma}

\begin{proof}
The matrix $[\mu_X(\varphi f_ig_j)]$ equals $C_f[\mu_X(\varphi u^{r+s})]_{r,s}C_g^{\mathsf T}$
by bilinearity. For the second statement, the matrix of $(\beta_i)$ is lower triangular with
diagonal $1/i!$.
\end{proof}

\emph{The expansion over summation points.} Let $\vvec(u)=(u^i)_{0\le i<h}$ be the column of monomials and $\vvec'(u)$ its derivative. For
$X=\xi$, \cref{lem:pole-series} and $(u^{i+j})'=(u^i)'u^j+u^i(u^j)'$ give
\begin{equation}\label{eq:real-sites}
  \bigl[\mu_\xi(\varphi u^{i+j})\bigr]_{i,j}
  =\sum_{m\ge1}w^m\,\bigl[\vvec(m)\ \vvec'(m)\bigr]
  \begin{pmatrix}\varphi'(m)&\varphi(m)\\ \varphi(m)&0\end{pmatrix}
  \bigl[\vvec(m)\ \vvec'(m)\bigr]^{\mathsf T},
\end{equation}
an absolutely convergent sum of matrices of rank at most two. We call the two columns
$\vvec(m)$ and $\vvec'(m)$ the \emph{slots} of the site $m$. The same structure, with other
sites and with additional rank-one terms at the poles, appears $2$-adically and $p$-adically in
\cref{sec:prime-two,sec:exceptional}. All of these expansions are used through the
Cauchy--Binet formula: if $M=EBE^{\mathsf T}$ with $E$ of size $h\times N$ and $B$ of size
$N\times N$, then
\begin{equation}\label{eq:cauchy-binet}
  \det M=\sum_{S,S'}\det E_{[h],S}\,\det B_{S,S'}\,\det E_{[h],S'},
\end{equation}
the sum running over pairs of $h$-element subsets of the column indices of $E$. When $B$ is
block diagonal, $\det B_{S,S'}=0$ unless $S$ and $S'$ meet every block in the same number of
indices, and then it is the product of the corresponding minors of the blocks.

\emph{Rank-one terms.} At a prime where the series over $m$ is not used, the pole part of the
matrix is kept as a sum of rank-one matrices, and the determinant is expanded as follows.

\begin{lemma}[rank-one expansion]\label{lem:rank-one}
Let $\mathbf P$ be an $h\times h$ matrix, $\mathbf c_0,\dots,\mathbf c_{K-1}$ columns of
length $h$ and $\gamma_0,\dots,\gamma_{K-1}$ scalars, all in a commutative ring. Then
\[
  \det\Bigl(\mathbf P+\sum_{k<K}\gamma_k\mathbf c_k\mathbf c_k^{\mathsf T}\Bigr)
  =\sum_{\mathcal K}\Bigl(\prod_{k\in\mathcal K}\gamma_k\Bigr)\,\Theta_{\mathcal K},
\]
where $\mathcal K$ runs over the subsets of $\{0,\dots,K-1\}$ with at most $h$ elements and
$\Theta_{\mathcal K}$ is a sum of products $\pm\det V_{I,\mathcal K}\det V_{J,\mathcal K}
\det\mathbf P_{I^c,J^c}$ over sets $I,J$ of $\abs{\mathcal K}$ rows, $V=[\mathbf c_k]_{k<K}$.
\end{lemma}

\begin{proof}
The matrix is $[I_h\mid V]\diag(\mathbf P,\Gamma)[I_h\mid V]^{\mathsf T}$ with
$\Gamma=\diag(\gamma_k)$. In~\eqref{eq:cauchy-binet} the middle minor vanishes unless both
index sets contain the same set $\mathcal K$ of columns of $V$, and then it is
$\prod_{k\in\mathcal K}\gamma_k$ times a minor of $\mathbf P$; the outer minors are, up to sign,
minors of $V$ with columns $\mathcal K$.
\end{proof}

\subsection{The criterion}\label{subsec:criterion}

\begin{proposition}\label{prop:criterion}
Let $\xi\in\R$ and let $(T_n)_{n\in\mathcal N}$, $\mathcal N\subset\N$ infinite, be nonzero
polynomials in $\Q[X]$ with $\deg T_n\le h_n$, $h_n\to\infty$. Suppose that, along $\mathcal N$,
\[
  \liminf\frac1{h_n^2}\sum_p\vG p(T_n)\log p\ \ge\ \aconst,\qquad
  \limsup\frac1{h_n^2}\log\abs{T_n(\xi)}\ \le\ \bconst,
\]
with $\bconst<\aconst$, and that for every $r\in\Q$ there are infinitely many $n\in\mathcal N$
with $T_n(r)\ne0$. Then $\xi\notin\Q$.
\end{proposition}

\begin{proof}
Let $\ell_n=\prod_pp^{-\vG p(T_n)}$; then $\ell_nT_n\in\Z[X]$. If $\xi=r/s$ with $s\ge1$, then
$s^{h_n}\ell_nT_n(r/s)$ is an integer, nonzero for infinitely many $n$, while
$\log\abs{s^{h_n}\ell_nT_n(\xi)}\le h_n\log s+(\bconst-\aconst+o(1))h_n^2\to-\infty$.
\end{proof}

This is Krisman's criterion \cite[Prop.~3.1, Thm.~3.2]{Krisman2026chi2} in the form we need;
its elementary core is formalised in Lean (see the section on reproducibility). We apply it with
$\mathcal N$ the set of all $n$ with $T_n\ne0$ for the arithmetic and analytic estimates, and with
the prime values of $n$ for nonvanishing.

\section{The prime 2}\label{sec:prime-two}

This section proves the lower bounds for $\vG2(T)$ that give the term $c_2\log2$ of
$\aconst_\sigma(a,q)$; they are used in \cref{thm:arithmetic}. The prime $2$ needs its own
arguments because the poles $x_k=-k-\frac12$ are not $2$-adically integral. There are two
cases. If $a$ is odd, $y$ is a $2$-adic unit and the difficulty is that $1-w$ is not a unit
(\cref{subsec:two-odd}). If $a$ is even, $w$ is divisible by $4$ and the series over the
summation points converges $2$-adically (\cref{subsec:two-even}). For a sequence $(c_l)_{l\ge0}$
we write $(\Delta c)_l=c_{l+1}-c_l$, so that $\Delta^Nc_0=\sum_{l}(-1)^{N-l}\binom Nlc_l$.

\subsection{Odd numerator: forward differences}\label{subsec:two-odd}

In this subsection $a$ is odd. Then $y=q/a\in\Z_2^\times$, and we put
\[
  Q=\sigma y^2=\frac1w,\qquad \nu=v_2(Q-1)=v_2(q^2-\sigma a^2).
\]
Since odd squares are congruent to $1$ modulo $8$, we have $\nu\ge3$ if $\sigma=1$ and $\nu=1$
if $\sigma=-1$; moreover $w$ is a unit and $v_2(1-w)=\nu$.

\emph{Heuristic.} Take as row polynomials the products of the first pole factors and as column
polynomials the products of the last ones. Then $\varphi f_ig_j$ has either few consecutive
poles and no polynomial part, or no pole at all. In the first case the entry is a forward
difference of pole values, which are linked by the recursion $T_{l+1}=QT_l-(2l+3)^{-2}$, and a
forward difference of order $N$ of such a sequence is divisible by about $2^{N}N!$. In the
second case the entry is a moment sum of a polynomial of small degree.

\begin{lemma}\label{lem:fd-squares}
Let $s$ be an odd integer and $g_l=(s+2l)^{-2}$ for $l\ge0$. Then
$v_2(\Delta^kg_0)\ge k+v_2(k!)$ for every $k\ge0$.
\end{lemma}

\begin{proof}
By induction on $k$, $\Delta^k\bigl[(s+2l)^{-1}\bigr]_{l=0}=(-2)^kk!\prod_{l=0}^k(s+2l)^{-1}$,
an identity of rational functions of $s$. Differentiating with respect to $s$ gives
\[
  \Delta^kg_0=(-2)^kk!\,\prod_{l=0}^k\frac1{s+2l}\,\sum_{l=0}^k\frac1{s+2l},
\]
and all the numbers $s+2l$ are odd.
\end{proof}

\begin{lemma}\label{lem:fd-recursion}
Let $s$ be odd, $g_l=(s+2l)^{-2}$, and let $(T_l)_{l\ge0}$ be polynomials in $\Z_2[X]$ with
$T_{l+1}=QT_l-g_l$. Then for $m\ge1$
\[
  \vG2(\Delta^mT_0)\ \ge\
  \begin{cases}m+v_2(m!)-2-\floor{\log_2m}&\text{if }\nu\ge2,\\ m-1&\text{if }\nu=1.\end{cases}
\]
\end{lemma}

\begin{proof}
From $\Delta T_l=(Q-1)T_l-g_l$ we get by induction
\[
  \Delta^mT_0=(Q-1)^mT_0-\sum_{k<m}(Q-1)^{m-1-k}\Delta^kg_0 .
\]
The first term has valuation at least $m\nu$, and by \cref{lem:fd-squares} the $k$-th term of
the sum has valuation at least $(m-1-k)\nu+k+v_2(k!)$. If $\nu=1$ both are at least $m-1$. If
$\nu\ge2$, then $m\nu\ge m+v_2(m!)$ and $(m-1-k)\nu+k+v_2(k!)\ge 2m-2-k+v_2(k!)$; since
$v_2(n!)=n-s_2(n)$ and $1\le s_2(n)\le1+\floor{\log_2n}$,
\[
  v_2(m!)-v_2(k!)=m-k-s_2(m)+s_2(k)\le m-k+\floor{\log_2m}\qquad(0\le k<m),
\]
which gives the claim.
\end{proof}

Let $f_i(u)=\prod_{k<i}(u-x_k)$ and $g_j(u)=\prod_{K-j\le k<K}(u-x_k)$ for $0\le i,j<h$. They
are monic of degrees $i$ and $j$, so $\Delta(X)=\det[\mu_X(\varphi f_ig_j)]$ by
\cref{lem:change-of-basis}. For a pair $(i,j)$ put $\pi=i+j+1-K$.

\begin{proposition}[entries]\label{prop:two-odd-entries}
Let $a$ be odd and $0\le i,j<h\le K$.
\begin{enumerate}[(a)]
\item If $\pi\le0$, put $N=-\pi$. Then $\vG2\bigl(\mu_X(\varphi f_ig_j)\bigr)\ge0$, and if
$\sigma=1$ and $N\ge2$ it is at least $N-1-\floor{\log_2N}$.
\item If $\pi\ge1$, then $\vG2\bigl(\mu_X(\varphi f_ig_j)\bigr)$ is at least
$-(\nu-1)\pi-2-\floor{\log_2\pi}$ if $\sigma=1$, and at least $-\pi$ if $\sigma=-1$.
\end{enumerate}
\end{proposition}

\begin{proof}
(a) Since $i,j<h\le K$ and $i\le K-j-1$, the factors of $f_i$ and $g_j$ cancel the first $i$
and the last $j$ factors of $D_K$, and $\varphi f_ig_j=u\prod_{l=0}^{N}(u-x_{i+l})^{-1}$. Its
polynomial part is constant (zero if $N\ge1$), so it does not contribute to $\mu_X$. Its residue
at $x_{i+l}$ is $x_{i+l}\prod_{l'\ne l}(l'-l)^{-1}=(-1)^lx_{i+l}/(l!\,(N-l)!)$. Put
\[
  T_l=\sigma^{i+l}y^{2(i+l)+1}(X-\tau_{i+l})\qquad(l\ge0).
\]
By~\eqref{eq:functional},
\begin{equation}\label{eq:entry-fd}
  \mu_X(\varphi f_ig_j)=-\frac4{N!}\sum_{l=0}^{N}(-1)^l\binom Nlx_{i+l}T_l
  =\frac{4(-1)^{N+1}}{N!}\,\Delta^N\bigl[x_{i+l}T_l\bigr]_{l=0}.
\end{equation}
The coefficients of $T_l$ are $\pm y^{2(i+l)+1}$ and
$\pm\sum_{j'\le i+l}\sigma^{j'}y^{2(i+l)-2j'}/(2j'+1)^2$, which lie in $\Z_2$, and
by~\eqref{eq:tau} $T_{l+1}=QT_l-(2i+2l+3)^{-2}$. So \cref{lem:fd-recursion} applies to
$(T_l)_{l\ge0}$ with $s=2i+3$ and to $(T_{l+1})_{l\ge0}$ with $s=2i+5$. Since $x_{i+l}$ is
affine in $l$ with $\Delta x=-1$, the Leibniz rule for differences gives, for $N\ge1$,
\[
  \Delta^N\bigl[x_{i+l}T_l\bigr]_{l=0}=x_i\,\Delta^NT_0-N\,\Delta^{N-1}T_1,
\]
and $v_2(x_i)=-1$. For $N=0$ the entry is $-4x_iT_0$, and for $N=1$ it is $4(x_i\Delta T_0-T_1)$;
both are integral. Let $N\ge2$. If $\nu\ge2$, \cref{lem:fd-recursion} and
$v_2(N)+v_2((N-1)!)=v_2(N!)$ show that both $x_i\Delta^NT_0$ and $N\Delta^{N-1}T_1$ have
valuation at least $N+v_2(N!)-3-\floor{\log_2N}$, and multiplying by $4/N!$ gives
$N-1-\floor{\log_2N}$. If $\nu=1$ they have valuation at least $N-2$ and $v_2(N)+N-2$, so the
entry has valuation at least $2-v_2(N!)+N-2=s_2(N)\ge0$.

(b) Now $K-j\le i$ and $\varphi f_ig_j=R(u)=u\prod_{k=K-j}^{i-1}(u-x_k)$ (the product is empty
if $\pi=1$), a polynomial of degree $\pi$, with $R(u)=2^{-\pi}H(2u)$ for a monic $H\in\Z[z]$. Hence
$R'(u)=2^{1-\pi}\sum_{d<\pi}c_d\,2^du^d$ with $c_d\in\Z$. Since $\Delta^b(u^d)(0)$ is $b!$ times
a Stirling number of the second kind and vanishes for $d<b$, we get
$v_2(\Delta^bR'(0))\ge1-\pi+b+v_2(b!)$ for $0\le b<\pi$. As $w$ is a unit and $v_2(1-w)=\nu$, we
have $v_2(\Omega_b(w))=-\nu(b+1)$, so
\[
  v_2\bigl(\mu_X(R)\bigr)=v_2\Bigl(\sum_{b<\pi}\Delta^bR'(0)\,\Omega_b(w)\Bigr)\ \ge\
  \min_{0\le b<\pi}\bigl[1-\pi-\nu+v_2(b!)-(\nu-1)b\bigr].
\]
If $\nu=1$ this is at least $-\pi$. If $\nu\ge2$, use $v_2(b!)\ge b-1-\floor{\log_2\pi}$ for
$1\le b<\pi$: the bracket is at least $-\pi-\nu-(\nu-2)b-\floor{\log_2\pi}$, which is smallest
at $b=\pi-1$, where it equals $-(\nu-1)\pi-2-\floor{\log_2\pi}$; the term $b=0$ is $1-\pi-\nu$,
which is not smaller.
\end{proof}

\begin{theorem}[prime $2$, odd numerator]\label{thm:two-odd}
Let $a$ and $q$ be odd, $K=\kappa h$ with $1\le\kappa\le2$, and $\nu=v_2(q^2-\sigma a^2)$. Then
\[
  \vG2(T)\ \ge\
  \begin{cases}\bigl[2(\kappa-1)-(\nu-2)\Dk\bigr]h^2-2h\bigl(1+\log_2(2h)\bigr)&\text{if }\sigma=1,\\[2pt]
  \bigl[\kappa-1-\Dk\bigr]h^2-h\log_2(2h)&\text{if }\sigma=-1.\end{cases}
\]
\end{theorem}

\begin{proof}
For a permutation $\varpi$ of $\{0,\dots,h-1\}$ put $\pi_i=i+\varpi(i)+1-K$. Then
$\sum_i\pi_i=h^2-hK$, and by convexity of $t\mapsto\pos t$,
\begin{equation}\label{eq:convexity-sum}
  \begin{aligned}
  \sum_i\pos{\pi_i}&\le\sum_i\tfrac12\bigl(\pos{2i+1-K}+\pos{2\varpi(i)+1-K}\bigr)\\
  &=\sum_{i<h}\pos{2i+1-K}\le\frac{(2h-K)^2}{4}=\Dk h^2 ;\end{aligned}
\end{equation}
the last sum is $\bigl(\frac{2h-K}2\bigr)^2$ if $K$ is even and $\frac{(2h-K)^2-1}4$ if $K$ is
odd. Let $\sigma=1$. By \cref{prop:two-odd-entries}, and since $N\le K\le2h$ and $\pi\le2h$,
every entry has Gauss valuation at least $-\pi-(\nu-2)\pos\pi-2-\log_2(2h)$. Summing over the
entries of a term of the Leibniz expansion gives
$\vG2(\Delta)\ge(\kappa-1)h^2-(\nu-2)\Dk h^2-h(2+\log_2(2h))$. Let $\sigma=-1$. Every entry has
Gauss valuation at least $-\pos\pi$, so $\vG2(\Delta)\ge-\Dk h^2$. Finally, by
$v_2(n!)=n-s_2(n)$ and $s_2(K)\le1+\log_2K$,
\begin{align*}
  v_2(\nrm)&=h\bigl(K-s_2(K)\bigr)-2\sum_{i<h}\bigl(i-s_2(i)\bigr)\\
  &\ge hK-h(h-1)-h(1+\log_2K)\ \ge\ (\kappa-1)h^2-h\log_2(2h).\qedhere
\end{align*}
\end{proof}

\begin{remark}\label{rem:two-odd-cancellation}
For $\sigma=1$ the bound decreases with $\nu$, which is unbounded as $q$ varies. In
$\aconst_\sigma(a,q)$ this is compensated exactly: the odd primes dividing $q^2-\sigma a^2$ cost
$\Dk\log\bigl((q^2-\sigma a^2)/2^\nu\bigr)$ per $h^2$ (\cref{prop:p-divides-Y}), so the total
contribution of $2$ and of these primes is
\[
  \bigl[2(\kappa-1)-(\nu-2)\Dk\bigr]\log2-\Dk\log\frac{q^2-a^2}{2^\nu}
  =\bigl[2(\kappa-1)+2\Dk\bigr]\log2-\Dk\log(q^2-a^2)
\]
for $\sigma=1$, and $[\kappa-1-\Dk]\log2-\Dk\log\frac{q^2+a^2}{2}=(\kappa-1)\log2-\Dk\log(q^2+a^2)$
for $\sigma=-1$. The error terms of \cref{thm:two-odd} do not depend on $\nu$. In exact
computations for $h\le36$ the bound for $\sigma=-1$ is far from the true valuation
(Appendix~\ref{app:consistency}); we have not tried to improve it.
\end{remark}

\subsection{Even numerator: parity minors}\label{subsec:two-even}

In this subsection $a=2^ea_0$ with $e\ge1$ and $a_0$ odd, so $v_2(y)=-e$ and $v_2(w)=2e$. In
the coordinate $z=2u$ the poles become the odd integers $z_k=-2k-1$, while the summation points
become the even integers $2m$, and the gain comes from treating the two parities separately.
For $e=1$ this is the argument of~\cite[\S3]{AxiomForge2026numeratorTwo}.

Put $P_K(z)=\prod_{k<K}(z-z_k)$ and $F(z)=z/P_K(z)$. Since $D_K(u)=2^{-K}P_K(2u)$,
\begin{equation}\label{eq:phi-F}
  \varphi(u)=2^{K-1}F(2u).
\end{equation}
On the space spanned by $\Q[z]$ and the fractions $1/(z-z_k)$ define the affine functional
\[
  \tilde\mu_X(G)=\mom\bigl(m\mapsto G'(2m)\bigr)\ (G\in\Q[z]),\qquad
  \tilde\mu_X\Bigl(\frac1{z-z_k}\Bigr)=-\sigma^ky^{2k+1}(X-\tau_k).
\]
If $R(u)=G(2u)$ then $R'(m)=2G'(2m)$, and $1/(u-x_k)=2/(2u-z_k)$; comparing
with~\eqref{eq:functional} gives $\mu_X(R)=2\tilde\mu_X(G)$ on both kinds of generators. With
$H(X)=[\tilde\mu_X(Fz^{i+j})]_{0\le i,j<h}$ and~\eqref{eq:phi-F}, the $(i,j)$ entry of the matrix
defining $\Delta$ is $2^{K-i-j}H_{ij}$, hence
\begin{equation}\label{eq:Delta-H}
  \Delta(X)=2^{Kh-h(h-1)}\det H(X).
\end{equation}

The series $\xi_2=\sum_{j\ge0}\sigma^jy^{-2j-1}/(2j+1)^2$ converges in $\Z_2$, because its
$j$-th term has valuation $e(2j+1)$. Repeating the computation in the proof of
\cref{lem:pole-series} in $\Q_2$, where $\sum_mw^m(\cdots)$ converges because $v_2(w)=2e>0$,
gives
\[
  \sum_{m\ge1}w^m\frac{\mathrm d}{\mathrm dz}\Bigl(\frac1{z-z_k}\Bigr)\Big|_{z=2m}
  =-\sigma^ky^{2k+1}(\xi_2-\tau_k).
\]
Hence $\tilde\mu_X(1/(z-z_k))=\sum_mw^m(\cdots)-\sigma^ky^{2k+1}(X-\xi_2)$. Writing
$Fz^s=Q_s+\sum_kr_kz_k^s/(z-z_k)$ with $Q_s$ a polynomial and $r_k=\Res_{z=z_k}F$, and using
\cref{lem:moment-sums}(a) $2$-adically for $Q_s$, we obtain the identity in $\Q_2[X]$
\begin{multline}\label{eq:H-two-adic}
  H(X)=\sum_{m\ge1}w^m\bigl[\vvec(2m)\ \vvec'(2m)\bigr]
  \begin{pmatrix}F'(2m)&F(2m)\\F(2m)&0\end{pmatrix}
  \bigl[\vvec(2m)\ \vvec'(2m)\bigr]^{\mathsf T}\\
  +(X-\xi_2)\sum_{k<K}\gamma_k\,\vvec(z_k)\vvec(z_k)^{\mathsf T},
\end{multline}
with $\gamma_k=-\sigma^ky^{2k+1}r_k$ and the series converging entrywise in $\Q_2[X]$ for the
Gauss norm. The weights have the following valuations. Since $z_k-z_l=2(l-k)$ and $z_k$ is odd,
\begin{equation}\label{eq:gamma-two}
  v_2(\gamma_k)=-e(2k+1)-(K-1)-v_2(k!)-v_2((K-1-k)!)\ \ge\ -e(2k+1)-2(K-1) .
\end{equation}
For $m\ge1$ the integer $P_K(2m)$ is odd, so $F(2m)\in\Z_2$, and
$F'(2m)=P_K(2m)^{-1}-F(2m)\sum_k(2m-z_k)^{-1}\in\Z_2$.

For $t\ge0$ put
\[
  \Pa(t)=\sum_{j<t}\bigl(j+v_2(j!)\bigr)=t(t-1)-\sum_{j<t}s_2(j)=t^2+O(t\log(t+2)).
\]

\begin{lemma}[parity minors]\label{lem:parity-minors}
Let $h\ge2$, $L=\floor{\log_2(h-1)}$, $\varepsilon\in\{0,1\}$, and let $J\subset\{0,\dots,h-1\}$
have $t$ elements. Consider a $t\times t$ matrix whose rows are indexed by $J$ and whose columns
are of the form $(\zeta^i)_{i\in J}$ or $(i\zeta^{i-1})_{i\in J}$ with
$\zeta\in\varepsilon+2\Z$, and let $d$ be the number of columns of the second kind. Then its
determinant has $2$-adic valuation at least $\Pa(t)-d(1+L)$.
\end{lemma}

\begin{proof}
Let $B_j(z)=\prod_{i<j}(z-\varepsilon-2i)$, a monic polynomial in $\Z[z]$ of degree $j$. Then
$z^i=\sum_{j\le i}c_{ij}B_j(z)$ with $c_{ij}\in\Z$ and $c_{ii}=1$. For $\zeta=2s+\varepsilon$,
$s\in\Z$, we have $B_j(\zeta)=2^jj!\binom sj$, which is divisible by $2^{j+v_2(j!)}$. For the
derivative, the identity
\begin{equation}\label{eq:binomial-derivative}
  \frac{\mathrm d}{\mathrm ds}\binom sj=\sum_{l=1}^{j}\frac{(-1)^{l-1}}{l}\binom s{j-l}
\end{equation}
(compare coefficients of $T^j$ in $\partial_s(1+T)^s=\log(1+T)(1+T)^s$) gives
$B_j'(\zeta)=2^{j-1}j!\,\partial_s\binom sj$, of valuation at least
$j+v_2(j!)-1-\floor{\log_2j}$. By the Cauchy--Binet formula for the product of the $t\times h$
matrix $(c_{ij})_{i\in J,\,j<h}$ and the $h\times t$ matrix of the values $B_j(\zeta)$ or
$B'_j(\zeta)$, the determinant is an integral combination of minors with rows
$j_1<\dots<j_t$ of the latter matrix. Since $j\mapsto j+v_2(j!)$ is increasing and
$j_r\ge r-1$, each such minor has valuation at least $\sum_{r<t}(r+v_2(r!))-d(1+L)$.
\end{proof}

\begin{corollary}\label{cor:parity-laplace}
Let $E$ be an $h\times h$ matrix whose columns are $P$ columns $\vvec(\zeta)$ with $\zeta$ odd
and $h-P$ columns $\vvec(\zeta)$ or $\vvec'(\zeta)$ with $\zeta$ even, $d$ of them derivatives.
Then $v_2(\det E)\ge\Pa(P)+\Pa(h-P)-d(1+L)$.
\end{corollary}

\begin{proof}
Laplace expansion along the odd columns writes $\det E$ as a signed sum, over $P$-element sets
$J$ of rows, of products of a $P\times P$ minor with odd nodes and rows $J$ and an
$(h-P)\times(h-P)$ minor with even nodes and the complementary rows. Apply
\cref{lem:parity-minors} to both.
\end{proof}

\begin{theorem}[prime $2$, even numerator]\label{thm:two-even}
Let $a=2^ea_0$ with $e\ge1$ and $a_0$ odd, $q$ odd, $\sigma=\pm1$, and $K\ge h\ge2$. Let
$\lambda_0\le\lambda_1\le\cdots$ be the numbers~\eqref{eq:gamma-two} in increasing order. Then
\begin{multline*}
  \vG2(T)\ \ge\ v_2(\nrm)+Kh-h(h-1)\\
  +\min_{0\le P\le h}\Bigl[2\Pa(P)+2\Pa(h-P)-(h-P)(1+L)+\sum_{i<P}\lambda_i
  +\sum_{i=1}^{h-P}2e\lceil i/2\rceil\Bigr].
\end{multline*}
If $K=\kappa h$ with $\kappa\ge1$ fixed, then
\begin{gather}
  \vG2(T)\ \ge\ C_e(\kappa)\,h^2-O(h\log h),\qquad C_e(\kappa)=\min_{0\le x\le1}f_e(x),
  \label{eq:two-asymptotic}\\
  f_e(x)=2(\kappa-1)-2(e+1)\kappa x+(e+2)x^2+\bigl(\tfrac e2+2\bigr)(1-x)^2.\notag
\end{gather}
The minimum is attained at $x_e=\min\bigl(1,\frac{2(e+1)\kappa+e+4}{3e+8}\bigr)$. For
$\kappa=5/4$: $C_2=-337/112$ at $x_2=27/28$, and $C_3=-9/2$ at $x_3=1$.
\end{theorem}

\begin{proof}
By~\eqref{eq:Delta-H}, $\vG2(T)=v_2(\nrm)+Kh-h(h-1)+\vG2(\det H)$. Truncate the sum over $m$
in~\eqref{eq:H-two-adic} at $m\le M$; the truncated matrices converge to $H$ in the Gauss norm,
so it suffices to bound the determinant of each truncation uniformly in $M$. Write the
truncation as $EBE^{\mathsf T}$, where the columns of $E$ are the slots $\vvec(2m)$,
$\vvec'(2m)$ ($1\le m\le M$) and $\vvec(z_k)$ ($k<K$), and $B$ is block diagonal with the
$2\times2$ blocks $w^m\bigl(\begin{smallmatrix}F'&F\\F&0\end{smallmatrix}\bigr)(2m)$ and the
$1\times1$ blocks $(X-\xi_2)\gamma_k$. By~\eqref{eq:cauchy-binet}, $\det H$ is a sum of terms
indexed by pairs $(S,S')$ of sets of $h$ slots that select the same poles and the same number
$n_m\in\{1,2\}$ of slots at each selected site $2m$.

Fix such a term with $P$ poles and $h-P$ even slots. The middle factor has Gauss valuation at
least $\sum_{k\in S}v_2(\gamma_k)+\sum_m2em\,n_m$, since a $1\times1$ minor of the block at $2m$
is $w^mF'(2m)$ or $w^mF(2m)$ and the $2\times2$ minor is $-w^{2m}F(2m)^2$, and
$\vG2(X-\xi_2)=0$. This is at least $\sum_{i<P}\lambda_i+\sum_{i\le h-P}2e\lceil i/2\rceil$,
because each site offers at most two slots. The derivative--derivative entry of each block
vanishes, so a site carrying one slot on each side contributes at most one derivative column to
the two outer factors together, and a site carrying two slots contributes two; the total number
of derivative columns on both sides is therefore at most $h-P$. By \cref{cor:parity-laplace} the
two outer factors together have valuation at least $2\Pa(P)+2\Pa(h-P)-(h-P)(1+L)$. This proves
the first inequality.

For the asymptotic form, $v_2(n!)=n-s_2(n)$ gives $v_2(\nrm)=Kh-h^2+O(h\log h)$ and
$\Pa(t)=t^2+O(h\log h)$ for $t\le h$. The $P$ smallest of the lower
bounds~\eqref{eq:gamma-two} are those with $k=K-1,\dots,K-P$, whose sum is
$-e(2KP-P^2)-2(K-1)P=-2(e+1)KP+eP^2+2P$, and $\sum_{i\le E}2e\lceil i/2\rceil\ge eE^2/2$.
Collecting terms with $P=xh$, $E=(1-x)h$ gives~\eqref{eq:two-asymptotic}. Finally
$f_e'(x)=(3e+8)x-2(e+1)\kappa-e-4$ is increasing, which gives $x_e$; for $\kappa=5/4$,
$x_2=\frac{27}{28}$ and $x_3=1$, and substituting gives $f_2(\frac{27}{28})=-\frac{337}{112}$
and $f_3(1)=-\frac92$.
\end{proof}

\begin{remark}\label{rem:two-even-comparison}
For $e=1$ and $\kappa=7/6$ one finds $C_1=-140/99$, the constant
of~\cite{AxiomForge2026numeratorTwo}. In general $f_e(1)=-e(2\kappa-1)$, and
$x_e=1$ if and only if $e\ge(2-\kappa)/(\kappa-1)$, that is, for $\kappa=5/4$, as soon as
$e\ge3$: then only the terms with $h$ poles matter and the bound is the one that the rank-one
argument of \cref{prop:p-divides-a} would give. For $a=8$, $q=12061$ the finite bound of
\cref{thm:two-even} equals the exact value of $\vG2(T)$ for every $h=4n\le48$, and for $a=4$,
$q=995$ it lies below the exact value by at most $12$ for $h\le56$
(Appendix~\ref{app:consistency}). This is consistency evidence only.
\end{remark}

\section{Odd primes dividing \texorpdfstring{$aq(q^2-\sigma a^2)$}{aq(q²-σa²)}}\label{sec:exceptional}

This section bounds $\vG p(T)$ for the finitely many odd primes dividing $aq(q^2-\sigma a^2)$,
which give the terms $-(2\kappa-1)\log a_0$, $\frac13\log q$ and $-\Dk\log(q^2-\sigma a^2)$ of
$\aconst_\sigma(a,q)$, and proves a crude uniform bound for the other odd primes, used for small
primes in \cref{sec:ordinary}. All of them rest on one lemma on the partial fractions of the
normalised weights (\cref{lem:newton-quotient}).

Throughout, $p$ is an odd prime. Since $w=\sigma a^2/q^2$ and $1-w=(q^2-\sigma a^2)/q^2$, the
odd primes fall into four classes: those dividing $a$ (where $\abs w_p<1$), those dividing $q$
(where $\abs w_p>1$), those dividing $q^2-\sigma a^2$ (where $w$ is a unit but $1-w$ is not),
and the others, called \emph{ordinary}, for which $w$ and $1-w$ are units. As $a$ and $q$ are
coprime, the first three classes are disjoint. We work with the binomial form
$T(X)=\det[\mu_X(\Phi\beta_i\beta_j)]$ of \cref{lem:change-of-basis}, with $\Phi$ as
in~\eqref{eq:Phi}.

\subsection{Newton quotients}\label{subsec:newton}

\begin{lemma}\label{lem:newton-quotient}
Let $p$ be odd and let $B\in\Q[u]$ have degree at most $D$ and satisfy $B(\Z)\subset\Z_p$. Then
\[
  \frac{K!\,B(u)}{D_K(u)}=Q(u)+\sum_{k<K}\frac{\rho_k}{u-x_k}
\]
with a polynomial $Q$ and
\begin{enumerate}[(a)]
\item $\rho_k=(-1)^kK\binom{K-1}{k}B(x_k)\in\Z_p$;
\item $Q'=0$ if $D\le K$; otherwise $\deg Q'\le D-K-1$ and $p^{\lambda}Q'(\Z)\subset\Z_p$
with $\lambda=\floor{\log_pD}+\floor{\log_p(D-K)}$.
\end{enumerate}
\end{lemma}

\begin{proof}
A polynomial with $p$-integral values on $\Z$ is a $\Z_p$-linear combination of the $\binom un$
and therefore maps $\Z_p$ into $\Z_p$; as $p$ is odd, $\frac12\in\Z_p$. Hence the Newton
coefficients $b_t$ in $B(u)=\sum_tb_t\binom{u+K-\frac12}{t}$, which are finite differences of
$B$ at $\frac12-K$, lie in $\Z_p$, and $B(x_k)\in\Z_p$. Part~(a) follows from
$D_K'(x_k)=\prod_{l\ne k}(l-k)=(-1)^kk!\,(K-1-k)!$. For $t\ge K$ the first $K$ factors of
$t!\binom{u+K-1/2}{t}$ are those of $D_K(u)$, so
\[
  \frac{K!}{D_K(u)}\binom{u+K-\frac12}{t}=\binom{t}{K}^{-1}\binom{u-\frac12}{t-K},
\]
while for $t<K$ the quotient is a proper rational function. Thus
\[
  Q(u)=\sum_{t=K}^{D}b_t\binom tK^{-1}\binom{u-1/2}{t-K}.
\]
By Kummer's theorem
$v_p\binom tK\le\floor{\log_pt}$, and
$\frac{\mathrm d}{\mathrm du}\binom{u-1/2}{s}=\sum_{l=1}^{s}\frac{(-1)^{l-1}}{l}\binom{u-1/2}{s-l}$
(the identity used in the proof of \cref{lem:parity-minors}); the values of $\binom{u-1/2}{r}$
at integers lie in $\Z_p$. This gives~(b).
\end{proof}

We apply the lemma to $B_{ij}(u)=u\,\beta_i(u)\beta_j(u)$, of degree $D_{ij}=i+j+1$, for which
$K!B_{ij}/D_K=\Phi\beta_i\beta_j$. Writing $Q_{ij}$ and $\rho_{ij,k}$ for the corresponding
quotient and residues,
\begin{equation}\label{eq:entry-decomposition}
  \begin{gathered}
  \mu_X(\Phi\beta_i\beta_j)=\mom(Q_{ij}')-4\sum_{k<K}\rho_{ij,k}\,\sigma^ky^{2k+1}(X-\tau_k),\\
  \rho_{ij,k}=(-1)^kK\tbinom{K-1}{k}x_k\beta_i(x_k)\beta_j(x_k).\end{gathered}
\end{equation}
For both ratios $D_{ij}\le2h-1<3h$ and $2K-1<3h$; we put $L_p=\floor{\log_p(3h)}$, so that
$p^{2L_p}Q_{ij}'(\Z)\subset\Z_p$.

\begin{proposition}[ordinary odd primes, uniform bound]\label{prop:uniform}
If $p\nmid 2aq(q^2-\sigma a^2)$, then $\vG p(T)\ge-2hL_p$.
\end{proposition}

\begin{proof}
Here $w$ and $1-w$ are $p$-adic units. We bound the two parts of~\eqref{eq:entry-decomposition}.
By \cref{lem:newton-quotient}(b) and \cref{lem:moment-sums}(c) applied to $p^{2L_p}Q_{ij}'$,
$\mom(Q_{ij}')\in p^{-2L_p}\Z_p$. Since $y$ is a unit, the polynomial $y^{2k+1}(X-\tau_k)$ has
coefficients $y^{2k+1}$ and $\sum_{j\le k}\sigma^jy^{2k-2j}/(2j+1)^2$, of valuation at least
$-2\floor{\log_p(2K-1)}\ge-2L_p$. With \cref{lem:newton-quotient}(a), every entry has Gauss
valuation at least $-2L_p$, and the Leibniz expansion gives the claim.
\end{proof}

\subsection{Primes dividing \texorpdfstring{$q^2-\sigma a^2$}{q²-σa²}}\label{subsec:p-divides-Y}

\begin{proposition}\label{prop:p-divides-Y}
Let $p$ be an odd prime dividing $q^2-\sigma a^2$ and $f=v_p(q^2-\sigma a^2)$. If $K=\kappa h$
with $1\le\kappa\le2$, then
\[
  \vG p(T)\ \ge\ -f\sum_{i<h}\pos{2i+1-K}-2hL_p\ \ge\ -f\,\Dk\,h^2-2hL_p .
\]
\end{proposition}

\begin{proof}
Now $w$ is a unit and $v_p(1-w)=f$, so $v_p(\Omega_b(w))=-f(b+1)$. The Newton expansion of
$Q_{ij}'$ involves only $b\le D_{ij}-K-1$, and its coefficients lie in $p^{-2L_p}\Z_p$, so
$v_p(\mom(Q_{ij}'))\ge-2L_p-f\pos{i+j+1-K}$; the pole terms are as in \cref{prop:uniform},
because $y$ is a unit. In the Leibniz expansion, for every permutation,
\eqref{eq:convexity-sum} bounds the sum of the $\pos{i+\varpi(i)+1-K}$ by
$\sum_{i<h}\pos{2i+1-K}\le\Dk h^2$.
\end{proof}

At $\kappa=1$ this is the bound $\frac14h^2\log(q^2-\sigma a^2)$ of
\cite[Lemma~5.1, Thm.~5.3]{Krisman2026chi2}; the decrease of the factor to $\frac{25}{144}$
or $\frac9{64}$ is the main gain of the ratios $\kappa>1$.

\subsection{Primes dividing \texorpdfstring{$q$}{q}: reflection}\label{subsec:p-divides-q}

Let $p\mid q$ and $f=v_p(q)\ge1$; then $v_p(y)=f$ and $v_p(w)=-2f$. The series
\[
  \xi'_p=\sum_{j\ge0}\frac{\sigma^jy^{2j+1}}{(2j+1)^2}
\]
converges in $\Z_p$: its $j$-th term has valuation $(2j+1)f-2v_p(2j+1)\ge(2j+1)-2\log_3(2j+1)>0$,
which tends to infinity.

\begin{lemma}[reflected functional]\label{lem:reflection}
For every $R\in\Rsp$,
\[
  \mu_X(R)=-\sum_{m\ge0}w^{-m}R'(-m)\;-\;4\,(X+\sigma\xi'_p)\sum_{k<K}\sigma^ky^{2k+1}
  \Res_{u=x_k}R,
\]
where the series converges in $\Q_p$.
\end{lemma}

\begin{proof}
For polynomials this is \cref{lem:moment-sums}(b), since $\abs{w^{-1}}_p<1$. For
$R=1/(u-x_k)$, $R'(-m)=-4/(2k+1-2m)^2$. The terms $0\le m\le k$ give, with $j=k-m$,
$4\sum_{j\le k}\sigma^{k-j}y^{2k-2j}/(2j+1)^2=4\sigma^ky^{2k+1}\tau_k$, and the terms $m>k$
give, with $j=m-k-1$, $4\sigma^{k+1}y^{2k+1}\xi'_p$. Hence
$-\sum_mw^{-m}R'(-m)=4\sigma^ky^{2k+1}(\tau_k+\sigma\xi'_p)$, and adding
$-4\sigma^ky^{2k+1}(X+\sigma\xi'_p)$ gives $-4\sigma^ky^{2k+1}(X-\tau_k)=\mu_X(R)$.
\end{proof}

\begin{proposition}\label{prop:p-divides-q}
Let $p\mid q$ and $f=v_p(q)$. Then $\vG p(T)\ge\frac f3h^2-O(h\log h)$.
\end{proposition}

\begin{proof}
Let $\bvec(u)=(\beta_i(u))_{i<h}$. By \cref{lem:reflection} applied to $\Phi\beta_i\beta_j$,
\begin{multline*}
  \bigl[\mu_X(\Phi\beta_i\beta_j)\bigr]=\sum_{m\ge0}w^{-m}\bigl[\bvec\ \bvec'\bigr](-m)
  \begin{pmatrix}-\Phi'(-m)&-\Phi(-m)\\-\Phi(-m)&0\end{pmatrix}
  \bigl[\bvec\ \bvec'\bigr](-m)^{\mathsf T}\\
  +\sum_{k<K}\eta_k\,\bvec(x_k)\bvec(x_k)^{\mathsf T},
\end{multline*}
with $\eta_k=-4(X+\sigma\xi'_p)\sigma^ky^{2k+1}\rho_k$ and
$\rho_k=\Res_{x_k}\Phi=(-1)^kK\binom{K-1}kx_k\in\Z_p$
(\cref{lem:newton-quotient}(a)). We record the valuations.
\begin{itemize}
\item $\bvec(-m)$ and $\bvec(x_k)$ are integral, and $\beta_i'=\sum_{l=1}^{i}\frac{(-1)^{l-1}}{l}
\beta_{i-l}$ shows that $\bvec'(-m)$ has valuation at least $-\floor{\log_ph}$.
\item $\vG p(\eta_k)\ge(2k+1)f$, because $X+\sigma\xi'_p$ has Gauss valuation $0$.
\item $v_p(w^{-m})=2fm$.
\item $2^KD_K(-m)=\prod_{k<K}(2k+1-2m)$ is a product of $K$ consecutive odd integers, none zero.
For each $e\ge1$ at most $\floor{K/p^e}+1$ of them are divisible by $p^e$, and none if
$p^e>2(K+m)$. Hence $v_p(K!/D_K(-m))\ge-\floor{\log_p(2K+2m)}$. Differentiating
$\Phi=K!\,u/D_K$ by the product rule, and using
$D_K'/D_K(-m)=\sum_k2/(2k+1-2m)$, we get $v_p(\Phi(-m))$, $v_p(\Phi'(-m))\ge
-C\log(h+m+2)$ for a constant $C$.
\end{itemize}
Truncate the sum at $m\le M$ and apply~\eqref{eq:cauchy-binet} with the slots $\bvec(-m)$,
$\bvec'(-m)$ and $\bvec(x_k)$. A nonzero term selects at each site $-m$ the same number of slots
on both sides and the same poles. Its outer factors have valuation at least $-2h\floor{\log_ph}$,
and its middle factor at least the sum, over the $h$ selected slots, of $2fm-C\log(h+m+2)$ for a
slot at $-m$ and of $(2k+1)f\ge2fk$ for the pole $x_k$. Every level $m\ge0$ thus offers at most
three slots, each of cost at least
\[
  2fm-C\log(h+m+2)\ \ge\ \bigl(2f-C/(h+2)\bigr)m-C\log(h+2),
\]
and the sum of the $h$ smallest such costs is at least
$(2f-C/(h+2))\sum_{i<h}\floor{i/3}-Ch\log(h+2)=\frac f3h^2-O(h\log h)$, uniformly in $M$. Since
$w^{-m}\to0$ in $\Q_p$, the truncated determinants converge to $T$ in the Gauss norm.
\end{proof}

This is Krisman's reflection argument \cite[Lemma~5.2]{Krisman2026chi2}, with prices $2fm$ at
the integer sites and $(2k+1)f$ at the poles; the numerator $a$, a $p$-adic unit, plays no role.

\subsection{Odd primes dividing the numerator}\label{subsec:p-divides-a}

Let $p$ be an odd prime dividing $a$ and $f=v_p(a)\ge1$; then $v_p(y)=-f$ and $v_p(w)=2f$.
\emph{Heuristic.} The pole $x_k$ enters with the factor $y^{2k+1}$, of valuation $-(2k+1)f$,
so a single entry of the matrix can have valuation about $-2Kf$. But the contribution of each
pole to the matrix has rank one, so a term of the determinant contains the factor $y^{2k+1}$ of
each pole at most once.

\begin{proposition}\label{prop:p-divides-a}
Let $p$ be an odd prime dividing $a$ and $f=v_p(a)$. Then
\[
  \vG p(T)\ \ge\ -f\,h(2K-h)-2hL_p=-f(2\kappa-1)h^2-2hL_p\qquad(K=\kappa h).
\]
\end{proposition}

\begin{proof}
By~\eqref{eq:entry-decomposition},
\[
  \bigl[\mu_X(\Phi\beta_i\beta_j)\bigr]=\mathbf P+\sum_{k<K}\gamma_k\,\bvec(x_k)\bvec(x_k)^{\mathsf T},
  \qquad\gamma_k=-4(-1)^kK\tbinom{K-1}kx_k\,\sigma^ky^{2k+1}(X-\tau_k),
\]
with $\mathbf P=[\mom(Q_{ij}')]$. We record the valuations.
\begin{itemize}
\item $\mathbf P$ has entries in $p^{-2L_p}\Z_p$: here $w\in p\Z_p$ and $1-w$ is a unit, so
\cref{lem:moment-sums}(c) applies to $p^{2L_p}Q_{ij}'$.
\item $\bvec(x_k)\in\Z_p^h$, because $x_k\in\Z_p$.
\item $\tau_k\in\Z_p$: its $j$-th term $\sigma^jy^{-2j-1}/(2j+1)^2$ has valuation
$(2j+1)f-2v_p(2j+1)>0$. Hence $\vG p(X-\tau_k)=0$ and $\vG p(\gamma_k)\ge-(2k+1)f$.
\end{itemize}
By \cref{lem:rank-one}, $T$ is a sum of terms $\prod_{k\in\mathcal K}\gamma_k\cdot\Theta_{\mathcal K}$
over sets $\mathcal K$ of at most $h$ poles, where $\Theta_{\mathcal K}$ is an integral
combination of products of two minors of the integral matrix $[\bvec(x_k)]_k$ and a minor of
size $h-\abs{\mathcal K}$ of $\mathbf P$. Hence
\[
  \vG p(T)\ \ge\ \min_{\mathcal K}\Bigl[-f\sum_{k\in\mathcal K}(2k+1)-2L_p(h-\abs{\mathcal K})\Bigr]
  \ \ge\ -f\sum_{k=K-h}^{K-1}(2k+1)-2hL_p,
\]
and $\sum_{k=K-h}^{K-1}(2k+1)=K^2-(K-h)^2=h(2K-h)$.
\end{proof}

\begin{remark}\label{rem:worse-sign}
Summing \cref{prop:p-divides-a} over the odd primes dividing $a$, \cref{prop:p-divides-q} over
the primes dividing $q$ and \cref{prop:p-divides-Y} over the odd primes dividing
$q^2-\sigma a^2$ gives the terms $-(2\kappa-1)\log a_0$, $\frac13\log q$ and, together with the
prime $2$ when $a$ is odd (\cref{rem:two-odd-cancellation}), $-\Dk\log(q^2-\sigma a^2)$ of
$\aconst_\sigma(a,q)$. When $a$ is even, $c_2$ does not depend on $\sigma$ and
$q^2-a^2<q^2+a^2$, so $\aconst_+(a,q)\ge\aconst_-(a,q)$: it suffices to compare the size exponent
with $\aconst_-(a,q)$. In exact computations the valuations at the primes dividing $a$ approach
$-f(2\kappa-1)h^2$ from above (Appendix~\ref{app:consistency}).
\end{remark}

\section{Ordinary primes}\label{sec:ordinary}

This section proves the lower bound $\Iprof_\kappa h^2$ for the contribution of the ordinary
primes $p\nmid2aq(q^2-\sigma a^2)$ to $\sum_p\vG p(T)\log p$ (\cref{thm:ordinary}); it is the
largest term of $\aconst_\sigma(a,q)$ and is used in \cref{thm:arithmetic}. Primes $p<\delta h$
are handled by the crude bound of \cref{prop:uniform} and primes $p>2K$ cost nothing
(\cref{lem:large-primes}). The work is in the range $p\asymp h$, where we bound $\vG p(T)$ by a
small optimisation problem over the residue classes modulo $p$ (\cref{subsec:two-bounds}), pass to
a limiting profile in $p/h$ (\cref{subsec:profile}), and integrate it against $\log p$. The numerator
$a$ plays no role here beyond the requirement that $w$, $1-w$ and $y$ are $p$-adic units; the
section is the argument of~\cite{AxiomForge2026numeratorTwo} for the two ratios $\kappa=7/6$ and
$\kappa=5/4$ (in the notation of that paper, without zero band).

Throughout, $p$ is ordinary, so $y$, $w$ and $1-w$ are $p$-adic units.

\begin{lemma}\label{lem:large-primes}
If $p>2K$ is ordinary, then $\vG p(T)\ge0$.
\end{lemma}

\begin{proof}
All factorials in $\nrm$ involve integers smaller than $p$, so $v_p(\nrm)=0$. The polynomial
$D_K$ is monic with $p$-integral coefficients, so the quotient of $u^{i+j+1}$ by $D_K$ has
$p$-integral coefficients and so has its derivative; by \cref{lem:moment-sums}(c) its moment sum
is $p$-integral. The residues $x_k^{i+j+1}/D_K'(x_k)$ are $p$-integral because
$D_K'(x_k)=\pm k!(K-1-k)!$, and $\tau_k$ involves only the odd numbers $2j+1<2K<p$.
\end{proof}

The interesting range is $p\asymp h$. Fix $\delta>0$. For $h$ large, every prime
$p\ge\delta h$ satisfies $p^2>2K$, and we assume this from now on. Then every difference
$l-k$ of pole indices and every odd number below $2K$ has $p$-adic valuation at most one.

\subsection{Residue classes and local functionals}\label{subsec:distribution}

\emph{Heuristic.} Modulo $p$ the poles fall into residue classes. In the coordinate $u=c+pz$
around a class $c$, the poles of that class are at distance $O(1)$ while all others are units,
so the weight $\varphi$ has a definite $p$-adic size on each class. Splitting the series over
the residue classes of $m$ (\cref{lem:distribution}) turns $\mu_X$ into a sum of local
functionals, one per class, which are integral (\cref{lem:local-integrality}).

For $c\in\{1,\dots,p\}$ let $\Pi_c=\{k<K:x_k\equiv c\pmod p\}$, $d_c=\abs{\Pi_c}$, and for
$k\in\Pi_c$ let $z_k=(x_k-c)/p\in\Z_p$; these are pairwise incongruent modulo $p$, because
$\abs{l-k}<K<p^2$. Let $\rest{\Z_p}$ be the ring of restricted power series
$\sum_jg_jz^j$ with $g_j\in\Z_p$, $g_j\to0$, and $\rest{\Q_p}=\rest{\Z_p}[1/p]$. Put $v=w^p$;
then $v\equiv w\pmod p$, so $v$ and $1-v$ are units. Let $M_r(v)=\sum_{n\ge0}n^rv^n\in\Z_p$ be
the rational moment ($M_0=1/(1-v)$, and $M_r=\sum_{i=1}^r S(r,i)\,i!\,v^i/(1-v)^{i+1}$ with
Stirling numbers $S(r,i)$), and define the continuous linear map
\[
  \mathsf M_v\colon\rest{\Q_p}\to\Q_p,\qquad \mathsf M_v\Bigl(\sum_jg_jz^j\Bigr)=\sum_jg_jM_j(v),
\]
which maps $\rest{\Z_p}$ into $\Z_p$. Let $V_c=\rest{\Q_p}\oplus\bigoplus_{k\in\Pi_c}\Q_p\cdot
(z-z_k)^{-1}$ and define $\ell_c\colon V_c\to\Q_p[X]$ by $\ell_c(G)=\mathsf M_v(G')$ for
$G\in\rest{\Q_p}$ and, for $k\in\Pi_c$,
\begin{equation}\label{eq:ell-pole}
  \ell_c\Bigl(\frac1{z-z_k}\Bigr)=\frac{p^2}{w^c}\Bigl[\mu_X\Bigl(\frac1{u-x_k}\Bigr)
  -\sum_{\substack{b=1\\b\ne c}}^{p}\frac{w^b}{p}\,
  \mathsf M_v\Bigl(\frac{\mathrm d}{\mathrm dz}\,\frac{1}{b+pz-x_k}\Bigr)\Bigr].
\end{equation}
For $R\in\Rsp$ let $R_c(z)=R(c+pz)$. If $R$ is a polynomial, $R_c\in\Q_p[z]$; if
$R=1/(u-x_k)$ then $R_c=p^{-1}(z-z_k)^{-1}$ for $k\in\Pi_c$, and otherwise $c-x_k$ is a unit and
$R_c\in\rest{\Z_p}$. Thus $R_c\in V_c$.

\begin{lemma}[distribution identity]\label{lem:distribution}
For every $R\in\Rsp$, $\displaystyle\mu_X(R)=\sum_{c=1}^{p}\frac{w^c}{p}\,\ell_c(R_c)$.
\end{lemma}

\begin{proof}
For a polynomial $g$ one has $\sum_{m\ge1}g(m)w^m=\sum_{c=1}^pw^c\sum_{n\ge0}g(c+pn)w^{pn}$ in
$\Q[[w]]$; both sides are rational functions of $w$, so the identity holds at our $w$, in
$\Q_p$. Apply it to $g=R'$ for a polynomial $R$, using $(R_c)'(z)=pR'(c+pz)$. For
$R=1/(u-x_k)$ with $k\in\Pi_{c_0}$ the identity is equivalent to the definition~\eqref{eq:ell-pole}
of $\ell_{c_0}(1/(z-z_k))$.
\end{proof}

\begin{lemma}[integrality]\label{lem:local-integrality}
Let $c\in\{1,\dots,p\}$ and $P_c(z)=\prod_{k\in\Pi_c}(z-z_k)$.
\begin{enumerate}[(a)]
\item For $k\in\Pi_c$, $\ell_c(1/(z-z_k))\in\Z_p[X]$.
\item If $N\in\rest{\Z_p}$, then $G=N/P_c\in V_c$ and $\ell_c(G)\in\Z_p[X]$.
\end{enumerate}
\end{lemma}

\begin{proof}
(a) In $p^2\mu_X(1/(u-x_k))=-4\sigma^ky^{2k+1}(p^2X-p^2\tau_k)$ the number $p^2\tau_k$ is
$p$-integral, since $v_p(2j+1)\le1$ for $2j+1<2K$, and $y$ is a unit. In the sum over $b\ne c$,
$p^{-1}\frac{\mathrm d}{\mathrm dz}(b+pz-x_k)^{-1}=-(b+pz-x_k)^{-2}$ lies in $\rest{\Z_p}$.
(b) Division of $z^j$ by the monic polynomial $P_c\in\Z_p[z]$ has integral quotient and
remainder, so $N=QP_c+N_0$ with $Q\in\rest{\Z_p}$ and $N_0\in\Z_p[z]$ of degree less than $d_c$.
Then $G=Q+\sum_{k\in\Pi_c}\frac{N_0(z_k)}{P_c'(z_k)}\frac1{z-z_k}$, and $P_c'(z_k)$ is a unit
because the $z_k$ are pairwise incongruent. Apply~(a).
\end{proof}

\emph{Local weights.} For $k\notin\Pi_c$ the factor $c+pz-x_k$ is a unit of $\rest{\Z_p}$. Hence
\begin{equation}\label{eq:local-weight}
  \varphi(c+pz)=p^{-d_c}\,\frac{N_c(z)}{P_c(z)},\qquad
  N_c(z)=(c+pz)\,E_c(z),\quad E_c\in\rest{\Z_p}^\times,
\end{equation}
so $N_c\in\rest{\Z_p}$. We call $\omega_c=-d_c-1$ the \emph{base cost} of the class $c$.

\subsection{Two lower bounds}\label{subsec:two-bounds}

\begin{proposition}[jet bound]\label{prop:jet}
$\vG p(\Delta)\ge\mathcal J_p$, where $\mathcal J_p$ is the minimum of
$\sum_c\sum_{j<t_c}(\omega_c+2j)$ over all $(t_c)_{c}$ with $t_c\ge0$ and $\sum_ct_c=h$.
\end{proposition}

\begin{proof}
For polynomials $f$ of degree less than $h$ write $f(c+pz)=\sum_{i<h}f^{(c)}_ip^iz^i$; for
$f=u^r$ the Taylor coefficients $f^{(c)}_i=\binom ri c^{r-i}$ are integers. By
\cref{lem:distribution} and~\eqref{eq:local-weight},
\[
  \mu_X(\varphi u^ru^s)=\sum_{c=1}^p\sum_{i,j<h}(u^r)^{(c)}_i(u^s)^{(c)}_j\,\Lambda^{(c)}_{ij},
  \qquad
  \Lambda^{(c)}_{ij}=w^cp^{\,i+j-1-d_c}\,\ell_c\Bigl(\frac{N_c(z)z^{i+j}}{P_c(z)}\Bigr),
\]
and $\vG p(\Lambda^{(c)}_{ij})\ge i+j+\omega_c$ by \cref{lem:local-integrality}(b). So the matrix of
$\Delta$ is $\mathbf J^{\mathsf T}\Lambda\mathbf J$ with $\mathbf J$ integral (rows indexed by
pairs $(c,i)$) and $\Lambda=\diag(\Lambda^{(c)})$ block diagonal. In~\eqref{eq:cauchy-binet} a
nonzero term selects $t_c$ rows and $t_c$ columns in each block, with distinct orders $i$ and
distinct orders $j$; each term of the Leibniz expansion of such a minor has valuation at least
$t_c\omega_c+2\binom{t_c}2=\sum_{j<t_c}(\omega_c+2j)$.
\end{proof}

Let $\epsilon_c=1$ if $\Pi_c$ contains an index $k$ with $2k+1<p$, and $\epsilon_c=0$
otherwise. Two elements of $\Pi_c$ differ by a multiple of $p$, so at most one index of $\Pi_c$
is below $(p-1)/2$, and it is the smallest one.

\begin{proposition}[pole bound]\label{prop:raw}
$\vG p(\Delta)\ge\mathcal R_p$, where $\mathcal R_p$ is the minimum of
$\sum_c\bigl[\sum_{j<t_c}(\omega_c+2j)+2\epsilon_c\one_{t_c=d_c}\bigr]$ over all $(t_c)$ with
$0\le t_c\le d_c$ and $\sum_ct_c\le h$.
\end{proposition}

\begin{proof}
Write $\varphi u^s=Q_s+\sum_kx_k^s\rho_k/(u-x_k)$, $\rho_k=x_k/D_K'(x_k)$. The quotient
$Q_s$ of $u^{s+1}$ by the monic $D_K\in\Z_p[u]$ is integral, so the matrix $\mathbf P=
[\mu_X(Q_{i+j})]$ is integral, and the matrix of $\Delta$ is
$\mathbf P+\sum_k\gamma_k\vvec(x_k)\vvec(x_k)^{\mathsf T}$ with
$\gamma_k=\rho_k\mu_X(1/(u-x_k))$. For $k\in\Pi_c$: exactly $d_c-1$
factors $l-k$ of $D_K'(x_k)$ are divisible by $p$, each exactly once, and $x_k$ is integral; and
$\mu_X(1/(u-x_k))$ has Gauss valuation at least $-2$, and at least $0$ if $2k+1<p$ (because then
$\tau_k$ has no denominator divisible by $p$). Hence $\vG p(\gamma_k)\ge\omega_c+2\cdot\one_{2k+1<p}$.

Write the matrix as $[I_h\mid V]\diag(\mathbf P,\Gamma)[I_h\mid V]^{\mathsf T}$ with
$V=[\vvec(x_k)]_k$ and $\Gamma=\diag(\gamma_k)$. A nonzero term of~\eqref{eq:cauchy-binet}
selects the same set $\mathcal K$ of poles on both sides, $\abs{\mathcal K}\le h$, and some columns of
$I_h$. Its middle factor has valuation at least $\sum_{k\in\mathcal K}\vG p(\gamma_k)$. Each outer
factor is, up to sign, a minor of $V$ with columns $\mathcal K$; as a polynomial in the nodes it is
alternating with integer coefficients, hence divisible by $\prod_{k<l\in\mathcal K}(x_l-x_k)$, and
the $t_c=\abs{\mathcal K\cap\Pi_c}$ nodes of each class, whose differences are divisible by $p$, contribute
valuation $\binom{t_c}{2}$ to each outer factor. If
$t_c=d_c$, the index with $2k+1<p$, when it exists, is selected.
\end{proof}

The normalising factor contributes, by Legendre's formula and $p^2>2K>h$,
\begin{equation}\label{eq:vp-S}
  v_p(\nrm)=h\floor{\tfrac Kp}-2\sum_{i<h}\floor{\tfrac ip}
  =h\Bigl(\floor{\tfrac Kp}-2Q_p+\tfrac phQ_p(Q_p+1)\Bigr),\qquad Q_p=\floor{h/p}.
\end{equation}
Hence, for ordinary $p\ge\delta h$ and $h$ large,
\begin{equation}\label{eq:vp-T-combined}
  \vG p(T)\ \ge\ \max\bigl(-2h,\ v_p(\nrm)+\max(\mathcal J_p,\mathcal R_p)\bigr),
\end{equation}
the term $-2h$ being \cref{prop:uniform} (for $p\ge\delta h$ one has $L_p=\floor{\log_p3h}\le1$ once
$p>\sqrt{3h}$).

\subsection{The profile}\label{subsec:profile}

\emph{Heuristic.} For $p=xh$ the residue classes come in two types, according to whether they
contain $\floor{\kappa/x}$ or one more pole, and, within each, to whether the first pole is
improved; the number of classes of each type is $h$ times the length of an interval. The minima
$\mathcal J_p$ and $\mathcal R_p$ are then $h$ times the values of two filling problems whose
capacities are these lengths. We use the filling problems in their dual form, which makes both the
comparison with $\mathcal J_p$, $\mathcal R_p$ and the convexity in $x$ immediate.

\begin{lemma}[filling problems]\label{lem:filling}
Let $(e_i)_{i\in I}$ be real costs with capacities $m_i\ge0$, finitely many below any given
bound, and let $M>0$. For every selection $0\le y_i\le m_i$ with $\sum_iy_i=M$ and every
$\lambda\in\R$,
\[
  \sum_iy_ie_i\ \ge\ \lambda M-\sum_im_i\pos{\lambda-e_i},
\]
and the same holds for $\sum_iy_i\le M$ if $\lambda\le0$. Equality holds for the greedy selection
(cheapest costs first) and $\lambda$ equal to the largest cost it uses (respectively the minimum of
that cost and $0$).
\end{lemma}

\begin{proof}
$\sum_iy_ie_i=\lambda\sum_iy_i-\sum_iy_i(\lambda-e_i)$ and
$y_i(\lambda-e_i)\le m_i\pos{\lambda-e_i}$; if $\sum y_i\le M$ and $\lambda\le0$, then
$\lambda\sum_iy_i\ge\lambda M$. For the greedy selection every $i$ with $e_i<\lambda$ is used fully
and every $i$ with $e_i>\lambda$ not at all.
\end{proof}

\emph{The model.} Let $\kappa=K/h$ be fixed and let $x>0$. Put $d_0=\floor{\kappa/x}$,
$r_\kappa=\kappa-d_0x$, and for $t\in[0,x)$
\[
  d(t)=d_0+\one_{t<r_\kappa},\qquad
  \epsilon(t)=\one_{t<x/2}\,\one_{d(t)>0},\qquad \omega(t)=-d(t)-1.
\]
Here $t$ stands for $k_0/h$, where $k_0$ is the least pole index in a class (see the proof of
\cref{prop:profile}). For $\lambda\in\R$ put
\[
  F_\lambda(\omega)=\sum_{j\ge0}\pos{\lambda-\omega-2j},\qquad
  \widetilde F_\lambda(\omega,d,\epsilon)=\sum_{0\le j<d}\pos{\lambda-\omega-2j-2\epsilon\one_{j=d-1}},
\]
and define
\begin{align*}
  J(x)&=\sup_{\lambda\in\R}\Bigl[\lambda-\int_0^xF_\lambda(\omega(t))\diff t\Bigr],\\
  R(x)&=\sup_{\lambda\le0}\Bigl[\lambda-\int_0^x\widetilde F_\lambda\bigl(\omega(t),d(t),
  \epsilon(t)\bigr)\diff t\Bigr],\\
  N(x)&=\floor{\kappa/x}-2\floor{1/x}+x\floor{1/x}\bigl(\floor{1/x}+1\bigr),\\
  B(x)&=\max\bigl(-2,\ N(x)+\max(J(x),R(x))\bigr).
\end{align*}
By \cref{lem:filling}, $J(x)$ is the least cost of filling mass $1$ with capacity $\diff t$ at each
cost $\omega(t)+2j$ (the \emph{jet slots}), and $R(x)$ the least cost of filling mass at most $1$
with capacity $\diff t$ at each cost $\omega(t)+2j+2\epsilon(t)\one_{j=d(t)-1}$, $j<d(t)$ (the
\emph{pole slots}); this primal form is the one evaluated in Appendix~\ref{app:profile-certificate}.

\begin{proposition}\label{prop:profile}
For every $\delta>0$ there is $C_\delta$ such that for all large $h$ in the construction and all
ordinary primes $\delta h\le p\le2K$, $\vG p(T)\ge h\,B(p/h)-C_\delta$.
\end{proposition}

\begin{proof}
Let $x=p/h\in[\delta,2\kappa]$. By~\eqref{eq:vp-S}, $v_p(\nrm)=hN(x)$ exactly, so
by~\eqref{eq:vp-T-combined} it suffices to show that the minima $\mathcal J_p$ and $\mathcal R_p$
of \cref{prop:jet,prop:raw} satisfy $\mathcal J_p\ge hJ(x)-C_\delta$ and
$\mathcal R_p\ge hR(x)-C_\delta$.

\emph{Exact class data.} The poles $x_0,\dots,x_{p-1}$ lie in distinct classes, so the classes are
indexed by $k_0\in\{0,\dots,p-1\}$, the least pole index in the class. Since $K/p=\kappa/x$ and
$K-p\floor{K/p}=hr_\kappa$, the class $k_0$ has exactly $d(k_0/h)$ poles. Its first pole has
$2k_0+1<p$ if and only if $k_0\le(p-3)/2$, which agrees with $\one_{k_0/h<x/2}$ except for
$k_0=(p-1)/2$. Hence at most one class has data $(d_c,\epsilon_c)$ different from
$(d,\epsilon)(k_0/h)$.

\emph{Counting.} The model data are constant on each of the at most three intervals into which
$x/2$ and $r_\kappa$ cut $[0,x)$, and an interval $[u,v)$ contains between $h(v-u)-1$ and
$h(v-u)+1$ of the points $k_0/h$. Hence, for every $\lambda$,
\[
  \sum_cF_\lambda(\omega_c)\le h\int_0^xF_\lambda(\omega(t))\diff t+4\max_\omega F_\lambda(\omega),
\]
the maximum over $\omega\in[-d_0-2,\,-1]$, the range of all base costs. The same count applies to
$\sum_c\widetilde F_\lambda(\omega_c,d_c,\epsilon_c)$ and $\int_0^x\widetilde F_\lambda\diff t$, since the summand
depends only on the data $(d,\epsilon)(k_0/h)$, which are constant on the same intervals.

\emph{Comparison.} $\mathcal J_p$ is the least sum of $h$ elements of the multiset
$\{\omega_c+2j\}$, so \cref{lem:filling} gives $\mathcal J_p\ge h\lambda-\sum_cF_\lambda(\omega_c)$
for every $\lambda$. The supremum defining $J(x)$ is attained at the largest cost $\lambda^*$ used
by the greedy filling; as the levels $j<\lceil1/x\rceil$ have total capacity at least $1$,
$\lambda^*\le-1+2\lceil1/x\rceil$. With $\lambda=\lambda^*$ we obtain $\mathcal J_p\ge
hJ(x)-4\max_\omega F_{\lambda^*}(\omega)$, and $F_{\lambda^*}(\omega)\le(\lambda^*+d_0+4)^2/4$ is
bounded in terms of $\delta$, because $d_0\le\kappa/\delta$ and $1/x\le1/\delta$. In the same way
$\mathcal R_p$ is the least sum of at most $h$ elements of the multiset of pole slots,
\cref{lem:filling} applies with every $\lambda\le0$, so $\mathcal R_p\ge h\lambda-\sum_c\widetilde F_\lambda(\omega_c,d_c,\epsilon_c)\ge h\bigl[\lambda-\int_0^x\widetilde F_\lambda\diff t\bigr]-4(d_0+1)(d_0+2)$,
because $\widetilde F_\lambda\le(d_0+1)(d_0+2)$ for $\lambda\le0$; taking the supremum over $\lambda\le0$ gives
$\mathcal R_p\ge hR(x)-C_\delta$.
\end{proof}

\begin{lemma}[piecewise convexity]\label{lem:convexity}
Let $\mathfrak B$ be the set of points $1/j$, $\kappa/j$, $2\kappa/j$, $j\ge1$. On every open
interval of $(0,\infty)\setminus\mathfrak B$ the function $B$ is convex. Consequently
$\int_a^bB\ge(b-a)B(\frac{a+b}2)$ for consecutive points $a<b$ of $\mathfrak B\cup\{\delta\}$.
\end{lemma}

\begin{proof}
Let $U$ be an open interval disjoint from $\mathfrak B$. The integers $\floor{1/x}$ and $d_0$
change only at the points $1/j$ and $\kappa/j$, so they are constant on $U$; hence $N$ is
affine on $U$, and $x/2$, $x$ and $r_\kappa=\kappa-d_0x$ are affine. Two of the cut points $0$,
$x/2$, $r_\kappa$, $x$ coincide only if $r_\kappa=0$ ($x=\kappa/d_0$) or $r_\kappa=x/2$
($x=2\kappa/(2d_0+1)$); the remaining coincidences ($r_\kappa=x$, $x/2=0$) are impossible. Both
values lie in $\mathfrak B$, so on $U$ the cut points keep their order, the data $(d,\epsilon)$
are constant on each piece, and the lengths of the pieces are affine in $x$. For fixed $\lambda$ the
functions $x\mapsto\lambda-\int_0^xF_\lambda(\omega(t))\diff t$ and
$x\mapsto\lambda-\int_0^x\widetilde F_\lambda\diff t$ are therefore affine on $U$, so $J$ and $R$,
as suprema of affine functions, are convex, and so is $B$. A convex function on $(a,b)$ lies above
its supporting line at $\frac{a+b}2$, which gives the integral inequality.
\end{proof}

\subsection{The ordinary contribution}\label{subsec:ordinary-contribution}

\begin{theorem}\label{thm:ordinary}
For every $\delta>0$,
\[
  \liminf_{h\to\infty}\frac1{h^2}\sum_{p\ \mathrm{ordinary}}\vG p(T)\log p\ \ge\
  \int_\delta^{2\kappa}B(x)\diff x-2\delta .
\]
With $\delta=1/2000$ the right-hand side is at least $\Iprof_{7/6}=-2.113164137$ for
$\kappa=7/6$ and at least $\Iprof_{5/4}=-2.355639833$ for $\kappa=5/4$.
\end{theorem}

\begin{proof}
Split the ordinary primes into $p<\delta h$, $\delta h\le p\le2K$ and $p>2K$. The last range
contributes a nonnegative amount by \cref{lem:large-primes}. For the first, by
\cref{prop:uniform},
\[
  \sum_{p<\delta h}2hL_p\log p\le2h\,\theta(\delta h)+2h\sqrt{3h}\log(3h)=2\delta h^2+o(h^2),
\]
because $L_p=1$ for $\sqrt{3h}<p<\delta h$, with $\theta(x)=\sum_{p\le x}\log p\sim x$. In the middle
range, \cref{prop:profile} and $\sum_{p\le2K}\log p=O(h)$ reduce the claim to
$\sum_{\delta h\le p\le2K}B(p/h)\log p\ge h\int_\delta^{2\kappa}B-o(h)$. The function $B$ is bounded
on $[\delta,2\kappa]$ and continuous outside the finite set $\mathfrak B\cap[\delta,2\kappa]$,
hence Riemann integrable; approximating it from below by step functions and using
$\theta(\beta h)-\theta(\beta'h)=(\beta-\beta')h+o(h)$ for fixed $\beta'<\beta$ (the prime number
theorem) gives the claim. The finitely many excluded primes (those dividing $2aq(q^2-\sigma a^2)$)
do not affect the limit. The numerical values are lower bounds for the midpoint sums of
\cref{lem:convexity} over the $5999$ intervals cut out by $\mathfrak B$ in $[\delta,2\kappa]$,
for each of the two values of $\kappa$, computed exactly in rational arithmetic
(Appendix~\ref{app:profile-certificate}).
\end{proof}

\begin{remark}\label{rem:profile-k1}
For $\kappa=1$ the same computation gives $\int_\delta^2B-2\delta=-1.7169166\ldots
=\frac12-\frac{133}{60}-\frac\delta2$. Since $\log\nrm=\frac12h^2+O(h\log h)$ when $\kappa=1$,
this is Krisman's generic constant $\frac{133}{60}$ \cite[Thm.~4.5]{Krisman2026chi2} up to the
cut-off: at $\kappa=1$ the profile reproduces his class bound, and the gain for $\kappa>1$ does
not come from a sharper estimate at ordinary primes.
\end{remark}

\section{The size of \texorpdfstring{$T(\xi)$}{T(ξ)}}\label{sec:decay}

This section proves the upper bound $\limsup h^{-2}\log\abs{T(\xi)}\le\bconst$
(\cref{thm:decay}), the second input of the criterion; it is combined with the arithmetic bound in
\cref{sec:proofs}. \emph{Heuristic.} By~\eqref{eq:real-sites}, $T(\xi)$ is a sum over
configurations of $h$ points on the summation sites, at most two per site, each configuration
weighted by a product of Vandermonde factors and of the weights $w^m\varphi(m)$. This is a discrete
logarithmic gas in an external field, and its size is controlled by comparing the configurations
with a fixed step density $\rho$; the comparison costs only $O(h\log h)$ thanks to a discrete
energy inequality. The method is Krisman's \cite[\S6--7]{Krisman2026chi2}, in the form of
\cite{AxiomForge2026numeratorTwo} without zero band. The numerator $a$ enters only through
$W=\abs w=a^2/q^2$, so the bound is the same for both signs and for all numerators with the same
$W$.

Throughout, $K=\kappa h$ with $\kappa\in\{7/6,5/4\}$.

\subsection{Statement}\label{subsec:decay-statement}

A \emph{comparison density} is a function $\rho\colon\R\to[0,2]$, constant on the cells of a
finite partition of $[0,1]$ into intervals, vanishing outside $[0,1]$, with $\int\rho=1$. Put
\begin{align*}
  U(x)&=\int\rho(t)\log\abs{x-t}\diff t,\qquad E(\rho)=\int U(x)\rho(x)\diff x,\\
  V(x)&=x\log W+\psi(x)-\psi(x+\kappa)\qquad(x\ge0),\\
  G(x)&=2U(x)+V(x),\qquad C=\kappa\log\kappa+\tfrac32-\kappa,
\end{align*}
with $\psi(t)=t\log\abs t-t$, $\psi(0)=0$. Note that $\psi(x)-\psi(x+\kappa)=-\int_x^{x+\kappa}\log
t\diff t$: it comes from the poles.

\begin{theorem}\label{thm:decay}
Let $\rho$ be a comparison density and $\lambda\in\R$, and suppose that there is $\eta>0$ with
$G(x)-\lambda\le-\eta x$ for all $x\ge1$. Then
\[
  \limsup_{h\to\infty}\frac{1}{h^2}\log\abs{T(\xi)}\ \le\ \bconst(\rho,\lambda):=
  C-E(\rho)+\lambda+2\int_0^1\pos{G(x)-\lambda}\diff x .
\]
\end{theorem}

\begin{lemma}[tail condition]\label{lem:tail}
For $x\ge1$, $G(x)\le x\log W+(2-\kappa)\log x$. Consequently the hypothesis of \cref{thm:decay}
holds if
\[
  1+\log W<0\qquad\text{and}\qquad\log W<\lambda .
\]
\end{lemma}

\begin{proof}
For $x\ge1$ and $t\in[0,1]$, $\log\abs{x-t}\le\log x$, so $U(x)\le\log x$. Next,
$\int_x^{x+\kappa}\log t\diff t\ge\kappa\log x$. This gives the first inequality. Since
$0<2-\kappa\le1$ and $\log x\le x-1$, $G(x)\le x\log W+(x-1)$, so
$G(x)-\lambda\le-(\lambda-\log W)-(x-1)(-1-\log W)\le-\eta x$ with
$\eta=\min(\lambda-\log W,\,-1-\log W)>0$.
\end{proof}

\subsection{Three lemmas and the proof}\label{subsec:decay-lemmas}

\begin{lemma}[minors with derivative columns]\label{lem:derivative-minors}
Let $h\ge2$ and let $s_1,\dots,s_h$ be integers, each value occurring at most twice. Let
$\mathcal E$ be the $h\times h$ matrix whose $i$-th column is $\vvec(s_i)$ or $\vvec'(s_i)$, where
for a value occurring twice one column of each kind is taken. Then
\[
  \abs{\det\mathcal E}\le h^h\Bigl(1+\frac2h\Bigr)^{h^2}\prod_{i<j,\ s_i\ne s_j}\abs{s_i-s_j}.
\]
\end{lemma}

\begin{proof}
Write each derivative column as $\vvec'(s)=\frac1{2\pi i}\oint_{\abs{\zeta-s}=1/h}
\vvec(\zeta)(\zeta-s)^{-2}\diff\zeta$. By multilinearity $\det\mathcal E$ is an integral of
Vandermonde determinants $\prod_{i<j}(\zeta_j-\zeta_i)$, with $\zeta_i=s_i$ for value columns
and $\abs{\zeta_i-s_i}=1/h$ otherwise, against a kernel of total mass at most $h$ per derivative
column. If $s_i\ne s_j$ then $\abs{\zeta_i-\zeta_j}\le\abs{s_i-s_j}+2/h\le(1+2/h)\abs{s_i-s_j}$;
if $s_i=s_j$ then $\abs{\zeta_i-\zeta_j}=1/h\le1$.
\end{proof}

\begin{lemma}[a discrete energy inequality]\label{lem:polya}
For every finitely supported real sequence $(\nu_s)_{s\in\Z}$ with $\sum_s\nu_s=0$,
\[
  \sum_{s\ne t}\nu_s\nu_t\log\abs{s-t}\ \le\ \log2\sum_s\nu_s^2 .
\]
\end{lemma}

\begin{proof}
We may assume the support lies in $\{0,\dots,M-1\}$ with $M\ge2$. Define $f$ on $\Z_{\ge0}$ by
$f(0)=\log(2M)$, $f(d)=\log(M/d)$ for $1\le d\le M$ and $f(d)=0$ for $d\ge M$. Its second
differences $f(d-1)-2f(d)+f(d+1)$ are $0$ at $d=1$, $\log\frac{d^2}{d^2-1}$ for $2\le d<M$,
$\log\frac M{M-1}$ at $d=M$, and $0$ beyond, so $f$ is convex and decreases to $0$, whence
$f(d)=\sum_{k\ge1}\bigl(f(k-1)-2f(k)+f(k+1)\bigr)\pos{k-d}$. The kernel
$\pos{k-\abs{s-t}}$ counts the windows of $k$ consecutive integers containing $s$ and $t$, so it is
positive semidefinite, and therefore $\sum_{s,t}\nu_s\nu_tf(\abs{s-t})\ge0$. Expanding,
this sum equals $\log2\sum_s\nu_s^2+\log M\,(\sum_s\nu_s)^2-\sum_{s\ne t}\nu_s\nu_t\log\abs{s-t}$.
\end{proof}

This is the Pólya-type inequality of \cite[Lemma~6.4]{Krisman2026chi2}.

\begin{lemma}[discretisation]\label{lem:discretisation}
Let $\rho$ be a comparison density, $m_s=h\int_{(s-1)/h}^{s/h}\rho$ for $s\in\Z$, and
\[
  P_m(s)=\sum_{t\ne s}m_t\log\abs{s-t},\qquad E_m=\sum_{s\ne t}m_sm_t\log\abs{s-t}.
\]
Then $0\le m_s\le2$, $\sum_sm_s=h$, $m_s=0$ unless $1\le s\le h$, and for every integer $s$,
\[
  P_m(s)\le hU(s/h)+h\log h+6+2\log h,\qquad E_m\ge h^2E(\rho)+h^2\log h-4h(1+\log h).
\]
\end{lemma}

\begin{proof}
Let $\varrho(u)=\rho(u/h)$, which has mass $m_t$ on $I_t=[t-1,t]$, and note
$hU(s/h)=\int\varrho(u)\log\abs{s-u}\diff u-h\log h$ and
$h^2E(\rho)+h^2\log h=\iint\varrho(u)\varrho(u')\log\abs{u-u'}\diff u\diff u'$. For the first
inequality compare, cell by cell, $m_t\log\abs{s-t}$ with $\int_{I_t}\varrho(u)\log\abs{s-u}\diff
u$. For $t<s$ the discrete term is the smaller one. For $t=s$ the discrete term is absent and
the integral is at least $-2\int_0^1\abs{\log v}\diff v=-2$. For $t=s+e$, $e\ge1$, the excess is at
most $2\int_{e-1}^{e}(\log e-\log v)\diff v$, which is $2$ for $e=1$ and at most $2/(e-1)$ for
$e\ge2$; summing over $e\le h$ gives at most $4+2\log h$. For the second inequality, cells with
$t=t'$ contribute a nonpositive amount to the double integral, and for $\abs{t-t'}=e\ge1$ we
have $\log\abs{u-u'}\le\log(e+1)\le\log e+1/e$, so the double integral is at most
$E_m+\sum_tm_t\sum_{1\le e\le h}2\cdot2/e$.
\end{proof}

\begin{proof}[Proof of \cref{thm:decay}]
By~\eqref{eq:real-sites}, $T(\xi)=\nrm\lim_{M\to\infty}\det H^{(M)}$, where $H^{(M)}$ is the sum
over $1\le m\le M$. Apply~\eqref{eq:cauchy-binet}. A nonzero term selects, at each site $s$, the
same number $n_s\in\{0,1,2\}$ of slots on both sides, with $\sum_sn_s=h$; each multiplicity
vector $n=(n_s)$ arises from at most $4^h$ pairs of slot sets. For $s\ge1$,
$\varphi'/\varphi(s)=1/s-\sum_k1/(s-x_k)$ has absolute value at most $1+K\le3h$, so every minor of
the middle factor is bounded by $\prod_sb(s)^{n_s}$ with $b(s)=3hW^s\abs{\varphi(s)}$. With
\cref{lem:derivative-minors},
\begin{equation}\label{eq:decay-step1}
  \abs{\det H^{(M)}}\le4^hh^{2h}e^{4h}\sum_n\prod_sb(s)^{n_s}\prod_{s<t}\abs{s-t}^{2n_sn_t}.
\end{equation}

\emph{Energy.} Let $m$ be as in \cref{lem:discretisation} and $\nu=n-m$. Then $\sum_s\nu_s=0$ and
$\sum_s\nu_s^2\le2\sum_s(n_s^2+m_s^2)\le8h$. Expanding $n=m+\nu$ and using
$\sum_{s\ne t}\nu_sm_t\log\abs{s-t}=\sum_sn_sP_m(s)-E_m$, \cref{lem:polya} and
\cref{lem:discretisation} give
\begin{align*}
  \sum_{s\ne t}n_sn_t\log\abs{s-t}
  &=-E_m+2\sum_sn_sP_m(s)+\sum_{s\ne t}\nu_s\nu_t\log\abs{s-t}\\
  &\le h^2\log h-h^2E(\rho)+2h\sum_sn_sU(s/h)+O(h\log h).
\end{align*}

\emph{Weights.} For $s\ge1$, by concavity
$\log D_K(s)=\sum_{k<K}\log(s+k+\frac12)\ge\int_0^K\log(s+t)\diff t=\psi(s+K)-\psi(s)$. Since
$\psi(hx)=h\psi(x)+hx\log h$, with $x=s/h$,
\[
  \log b(s)\le hV(x)-K\log h+\log s+\log(3h).
\]

\emph{Summation.} Combining, the logarithm of the summand in~\eqref{eq:decay-step1} is at most
\[
  (1-\kappa)h^2\log h-h^2E(\rho)+h^2\lambda+\sum_sn_s\bigl(h(G(s/h)-\lambda)+\log s\bigr)+O(h\log h).
\]
For $s\le h$ we use $n_sh(G-\lambda)\le2h\pos{G-\lambda}$ and $\log s\le\log h$; for $s>h$ the
tail hypothesis gives $h(G(s/h)-\lambda)\le-\eta s$. Since every multiplicity vector comes from at
least one ordered $h$-tuple of sites,
\[
  \sum_n\prod_s\bigl(se^{h(G(s/h)-\lambda)}\bigr)^{n_s}\le
  \exp\Bigl(2h\sum_{1\le s\le h}\pos{G(s/h)-\lambda}\Bigr)\Bigl(h^2+\sum_{s>h}se^{-\eta s}\Bigr)^h,
\]
uniformly in $M$. Finally Stirling's formula gives
$\log\nrm=(\kappa-1)h^2\log h+Ch^2+O(h\log h)$, the terms $h^2\log h$ cancel, and since $G$ is
continuous on $[0,1]$ the Riemann sums $h^{-1}\sum_{1\le s\le h}\pos{G(s/h)-\lambda}$ converge
to $\int_0^1\pos{G-\lambda}$.
\end{proof}

\subsection{Larger denominators}\label{subsec:shift}

The exponent depends on $q$ only through $\log W=2\log a-2\log q$, and only in the term $x\log W$
of $V$.

\begin{proposition}[shift]\label{prop:shift}
Let $q_0\le q$ be odd, let $(\rho,\lambda)$ satisfy the conditions of \cref{lem:tail} for
$W_0=a^2/q_0^2$, and put $\tau=2\log(q/q_0)$. Then $(\rho,\lambda-\tau/2)$ satisfies them for
$W=a^2/q^2$, and $\bconst_q(\rho,\lambda-\tau/2)\le\bconst_{q_0}(\rho,\lambda)-\frac12\log(q/q_0)$.
\end{proposition}

\begin{proof}
We have $G_q(x)=G_{q_0}(x)-\tau x$, so $G_q-(\lambda-\tau/2)=G_{q_0}-\lambda-\tau(x-\frac12)$ and
$\pos{G_q-\lambda+\tau/2}\le\pos{G_{q_0}-\lambda}+\tau\pos{\frac12-x}$. Integrating over $[0,1]$
costs at most $2\cdot\tau/8$, while $\lambda$ decreases by $\tau/2$. The quantities
$\lambda-\log W$ and $-1-\log W$ of \cref{lem:tail} become $\lambda-\log W+\tau/2$ and $-1-\log W+\tau$, so they stay positive.
\end{proof}

This is Krisman's shift lemma \cite[Lemma~7.1]{Krisman2026chi2}. The bound of
\cref{thm:decay} depends on $\abs w$ only, so it is the same for $\chi_2$ and $\Ti_2$; Krisman
observes \cite[Remark~6.7]{Krisman2026chi2} that for $w<0$ the true determinants are smaller,
which we do not use.

\section{Nonvanishing at prime indices}\label{sec:nonvanishing}

This section proves that $T_n(r)\ne0$ for every rational $r$ and every large prime $n$
(\cref{thm:nonvanishing} with \cref{prop:local-dets}), the third input of the criterion, used in
\cref{sec:proofs}. \emph{Heuristic.} When $h$ and $K$ are fixed multiples of a prime $p$, every
residue class modulo $p$ contains the same number of poles. In a basis of polynomials adapted to the
classes, the distribution identity of \cref{sec:ordinary} makes the scaled matrix block triangular
modulo $p$, and each diagonal block reduces to one of three fixed rational matrices that do not
depend on $p$. If their determinants are nonzero, $\Delta(r)$ is a $p$-adic unit times an explicit
power of $p$. The numerator $a$ enters only through $Q=\sigma y^2=\sigma q^2/a^2$.

Let $(t,\gamma)\in\{(6,7),(4,5)\}$, and for a prime $p$ consider the determinant with
\[
  h=tp,\qquad K=\gamma p
\]
(the construction with $n=p$). Put $Q=\sigma y^2$, so that $w=1/Q$.

\subsection{The local matrices and the statement}\label{subsec:local-matrices}

For $s\in\{0,1\}$ let $\mathcal V_s$ be the $\Q$-vector space spanned by $\Q[z]$ and the
fractions $1/(z+l+\frac12)$, $s\le l<s+\gamma$, and define $\ell_Q\colon\mathcal V_s\to\Q$ by
\begin{equation}\label{eq:ellQ}
  \ell_Q(g)=\sum_{r\ge0}[z^r]g'\cdot M_r(1/Q)\quad(g\in\Q[z]),\qquad
  \ell_Q\Bigl(\frac1{z+l+\frac12}\Bigr)=4\sum_{i<l}\frac{Q^{\,l-i}}{(2i+1)^2},
\end{equation}
where $M_r(v)=\sum_{n\ge0}n^rv^n$ is the rational function of \cref{subsec:distribution}. For
$(s,e)\in\{(0,0),(1,0),(1,1)\}$ let $L_{s,e}(Q)$ be the $t\times t$ matrix with entries
\[
  \ell_Q\Bigl(\frac{z^{i+j}(z+1)^e}{\prod_{l=s}^{s+\gamma-1}(z+l+\frac12)}\Bigr),\qquad 0\le i,j<t .
\]
Its entries are rational numbers depending on $Q$ and $(t,\gamma)$ only.

\begin{theorem}\label{thm:nonvanishing}
Assume that the three determinants $\det L_{s,e}(Q)$ are nonzero. Then for every $r\in\Q$ and
every sufficiently large prime $p$ one has $\Delta(r)\ne0$, hence $T(r)\ne0$. More precisely,
$v_p(\Delta(r))=-pt(\gamma+2-t)+t$.
\end{theorem}

For $(t,\gamma)=(6,7)$ the valuation is $-18p+6$ and for $(4,5)$ it is $-12p+4$. We checked this
prediction against exact computations of $\Delta(0)$ at primes $p=13,17$ for five of the six
threshold rows, with $h=52,68,102$ (Appendix~\ref{app:consistency}).

\begin{proof}[Proof of \cref{thm:nonvanishing}]
Let $p$ be a prime not dividing $2aq(q^2-\sigma a^2)$ or the denominator of $r$, with $p>2\gamma+1$
(so $p^2>2K$), and not dividing the numerators and denominators of the three determinants or the
denominators of the entries of the $L_{s,e}(Q)$. All but finitely many primes qualify. We use the
notation of \cref{sec:ordinary}; every residue class $c\in\{1,\dots,p\}$ contains exactly
$\gamma$ poles.

\emph{Step 1: a basis adapted to the classes.} The ideals $((u-c)^t)$, $1\le c\le p$, of
$\Z_p[u]$ are pairwise comaximal, because $c-c'$ is a unit for $c\ne c'$. By the Chinese remainder
theorem there are $f_{c,i}\in\Z_p[u]$ of degree less than $tp$, for $1\le c\le p$ and $0\le
i<t$, with $f_{c,i}\equiv(u-c)^i\pmod{(u-c)^t}$ and $f_{c,i}\equiv0\pmod{(u-c')^t}$ for $c'\ne c$.
They form a $\Z_p$-basis of the polynomials of degree less than $h$ with $\Z_p$ coefficients, so
their coefficient matrix $C$ has $\det C\in\Z_p^\times$, and by \cref{lem:change-of-basis} the
Gram matrix $\mathbf G=[\mu_r(\varphi f_{c,i}f_{c',j})]$ has $\det\mathbf G=(\det C)^2\Delta(r)$.
Write $f_{c,i}(b+pz)=p^iF_{c,i,b}(z)$; then $F_{c,i,b}\in\Z_p[z]$, $F_{c,i,c}\equiv z^i$ and,
for $b\ne c$, $F_{c,i,b}\in p^{t-i}\Z_p[z]\subset p\Z_p[z]$.

\emph{Step 2: local weights.} By~\eqref{eq:local-weight}, for $b\ne p$,
$\varphi(b+pz)=p^{-\gamma}\Psi_b(z)$ with $\Psi_b=N_b/P_b$, $N_b\in\rest{\Z_p}$ and
\[
  N_b\equiv\lambda_b\pmod p,\qquad P_b(z)=\prod_{l\in J_b}\bigl(z+l+\tfrac12\bigr),
\]
where $\lambda_b=b\prod_{k\notin\Pi_b}(b-x_k)^{-1}$ is a unit. Here $J_b=\{0,\dots,\gamma-1\}$ if
$b\le\frac{p-1}2$ and $J_b=\{1,\dots,\gamma\}$ if $b\ge\frac{p+1}2$: indeed
$x_k-b=-(2k+2b+1)/2$, and $2k+2b+1=(2l+1)p$ for $0\le k<\gamma p$ gives exactly these ranges of
$l$, with $z_k=-l-\frac12$ \emph{exactly}. For $b=p$ the numerator $u=p(1+z)$ supplies one more
factor $p$: $\varphi(p+pz)=p^{1-\gamma}\Psi_p(z)$ with $N_p\equiv\lambda_p(1+z)$, where now $\lambda_p=\prod_{k\notin\Pi_p}(-x_k)^{-1}$ is a unit.

\emph{Step 3: the scaled matrix.} By \cref{lem:distribution},
$\mathbf G_{(c,i),(c',j)}=\sum_bw^bp^{\,i+j-1-\gamma}\ell_b(\tilde\varphi_bF_{c,i,b}F_{c',j,b})$, where
$\tilde\varphi_b=\Psi_b$ for $b\ne p$ and $\tilde\varphi_p=p\Psi_p$. Let $\mathbf M$ be obtained from
$\mathbf G$ by multiplying row $(c,i)$ by $p^{\gamma+1-i}$, and moreover by $p^{-1}$ if $c=p$, and
column $(c',j)$ by $p^{-j}$. Then
\[
  \mathbf M_{(c,i),(c',j)}=p^{-\one_{c=p}}\sum_{b=1}^pw^b\,\ell_b\bigl(\tilde\varphi_bF_{c,i,b}F_{c',j,b}\bigr),
\]
which is $p$-integral by \cref{lem:local-integrality}(b): for $c\ne p$ directly; for $c=p$ because
$F_{p,i,b}\in p\Z_p[z]$ when $b\ne p$, and $\tilde\varphi_p\in p\,V_p$.

\emph{Step 4: block triangularity.} Modulo $p$, $F_{c,i,b}\equiv\one_{c=b}z^i$, and $\ell_b$ is
$\Z_p$-linear on the lattice of \cref{lem:local-integrality}(b), so a term with $b\ne c$ or $b\ne c'$
vanishes modulo $p$ unless it is divided by $p$. Hence for $c\ne p$ the row $(c,i)$ of
$\mathbf M$ is, modulo $p$, supported in the columns $(c,j)$, with entries
$w^c\ell_c(\Psi_cz^{i+j})$. For $c=p$, the entries in columns $(p,j)$ receive from $b\ne p$ terms
in $p^{2t-i-j-1}\Z_p\subset p\Z_p$, so they are $w^p\ell_p(\Psi_pz^{i+j})$ modulo $p$. Ordering the
classes with $p$ last, $\mathbf M$ is block lower triangular modulo $p$, and
\begin{equation}\label{eq:det-product}
  \det\mathbf M\equiv\prod_{c=1}^{p}\det\bigl[w^c\ell_c(\Psi_cz^{i+j})\bigr]_{i,j<t}\pmod p .
\end{equation}

\emph{Step 5: reduction of the local functionals.} We claim that for $N\in\rest{\Z_p}$,
$\ell_c(N/P_c)\equiv\ell_Q(\bar N/P_c)\pmod p$, where $\bar N$ is the polynomial obtained by
reducing the coefficients of $N$ and $\ell_Q$ is reduced modulo $p$. For the restricted-series
part this follows from $v=w^p\equiv w=1/Q\pmod p$ and the fact that $M_r$ has denominator
$(1-v)^{r+1}$, a unit. For a pole $k\in\Pi_c$ with $z_k=-l-\frac12$, the sum over $b\ne c$
in~\eqref{eq:ell-pole} is divisible by $p^2$, and $p^2X\equiv0$, so
$\ell_c(1/(z-z_k))\equiv4\sigma^ky^{2k+1}w^{-c}p^2\tau_k$. In $p^2\tau_k$ only the terms with
$p\mid2i'+1$ survive; writing $2i'+1=(2i+1)p$, they are those with $i<l$ (because
$2k+1=(2l+1)p-2c$ and $1\le c\le p$), and they contribute
$\sigma^{ip+(p-1)/2}y^{-(2i+1)p}/(2i+1)^2$. Using $\sigma^p=\sigma$, $y^p\equiv y$, $w^{-c}=\sigma^c
y^{2c}$, $y^{2k+2c+1}=y^{(2l+1)p}$ and $k+c+\frac{p-1}2=(l+1)p-1\equiv l\pmod2$, we obtain
\[
  \ell_c\Bigl(\frac{1}{z-z_k}\Bigr)\equiv4\sum_{i<l}\sigma^{l+i}y^{2l-2i}(2i+1)^{-2}
  =4\sum_{i<l}\frac{Q^{\,l-i}}{(2i+1)^2}\pmod p,
\]
which is~\eqref{eq:ellQ}. Here $y^p\equiv y$ means congruence in $\Z_p$ modulo $p$ after
reduction, $y=q/a$ being a $p$-adic unit, and the reduction of rational numbers to $\mathbb F_p$
is the one used throughout. Since the decomposition of $N/P_c$ into a restricted series and
principal parts has integral coefficients (\cref{lem:local-integrality}(b)) and commutes with
reduction, the claim follows.

\emph{Step 6: conclusion.} By Steps 2 and 5 the $c$-th block in~\eqref{eq:det-product} is
congruent to $w^c\lambda_c$ times $L_{s_c,e_c}(Q)$, with $s_c=\one_{c>(p-1)/2}$ and $e_c=\one_{c=p}$.
Hence $\det\mathbf M$ is a $p$-adic unit. As $\det\mathbf M=p^{\,pt(\gamma+1)-pt(t-1)-t}\det\mathbf
G$ and $\det C$ is a unit, $v_p(\Delta(r))=-pt(\gamma+2-t)+t$.
\end{proof}

Krisman's Lemma~P \cite[Thm.~8.1]{Krisman2026chi2} follows the same pattern with $h=3\ell$,
$K=h$, the Lagrange basis at the poles and three $3\times3$ blocks. With $K>h$ there is no
Lagrange basis at the poles, which is why we use the Taylor jets of Step~1.

\subsection{The local determinants}\label{subsec:local-dets}

\begin{proposition}\label{prop:local-dets}
For $(t,\gamma)=(6,7)$ and $(s,e)=(0,0),(1,0),(1,1)$,
\[
  (Q-1)^{24}\det L_{s,e}(Q)=c_{s,e}\,Q^{\,a_{s,e}}(Q-1)^{b_{s,e}}\,\Pi_{s,e}(Q)
\]
with nonzero rationals $c_{s,e}$, $(a_{s,e},b_{s,e})=(9,20),(11,20),(13,18)$, and primitive
integer polynomials $\Pi_{s,e}$ of degrees $16,20,20$. For $(t,\gamma)=(4,5)$ the same holds with
$(Q-1)^{16}$, $(a_{s,e},b_{s,e})=(5,15),(6,15),(7,14)$ and degrees $6,9,9$. In all six cases the
roots of $\Pi_{s,e}$ have modulus less than Cauchy's bound $1+\max_i\abs{\pi_i/\pi_{\deg}}$, which is less than
$23.19$, $136.73$, $188.09$ for $(6,7)$ and less than $1.71$, $9.38$, $5.08$ for $(4,5)$. In
particular $\det L_{s,e}(Q)\ne0$ whenever $\abs Q\ge189$, hence whenever $q\ge14a$.
\end{proposition}

\begin{proof}
The entries of $(Q-1)^4L_{s,e}(Q)$ are polynomials in $Q$: the derivative of the polynomial part has
degree at most $3$, so only $M_0,\dots,M_3$ occur, and the pole values are polynomials. The
determinants were computed exactly, twice and independently (Appendix~\ref{app:local-certificates}): once
symbolically in $Q$, and once by exact interpolation at rational values, with agreement. The bound on
the roots is Cauchy's, evaluated exactly. Finally $q\ge14a$ gives $\abs Q=q^2/a^2\ge196>189$.
\end{proof}

The proposition is used with $\abs Q\ge196$; all six threshold rows have $q/a\ge43$.

\section{Proof of the main theorem}\label{sec:proofs}

This section combines the three estimates through the criterion of \cref{prop:criterion}: it
collects the local bounds into the arithmetic constant (\cref{thm:arithmetic}), shows that the margin
between the arithmetic and the size constants grows with $q$ (\cref{prop:propagation}), compares them
at the six thresholds (\cref{thm:rows}), and records the status of every step of the proof
(\cref{subsec:step-status}).

\begin{theorem}[arithmetic]\label{thm:arithmetic}
Let $a\ge1$, $q>a$ odd and prime to $a$, $\sigma=\pm1$ and $\kappa\in\{7/6,5/4\}$. Along the $n$
with $T_n\ne0$, $\liminf h^{-2}\sum_p\vG p(T_n)\log p\ge\aconst_\sigma(a,q)$, with
$\aconst_\sigma(a,q)$ as in \cref{thm:three-estimates}.
\end{theorem}

\begin{proof}
Sum \cref{thm:two-odd} or \cref{thm:two-even} (prime $2$), \cref{prop:p-divides-a} over the odd
primes dividing $a$, \cref{prop:p-divides-q} over the primes dividing $q$,
\cref{prop:p-divides-Y} over the odd primes dividing $q^2-\sigma a^2$, and \cref{thm:ordinary}; these
sets of primes are disjoint and the first four are finite. Note $\sum_{p\mid a_0}v_p(a)\log
p=\log a_0$, $\sum_{p\mid q}v_p(q)\log p=\log q$, and, for odd $a$, the cancellation of
\cref{rem:two-odd-cancellation} between the prime $2$ and the term $-\Dk\log(q^2-\sigma a^2)$.
\end{proof}

\begin{proposition}[propagation in $q$]\label{prop:propagation}
Let $a$, $\sigma$, $\kappa$ be as above, let $q_0\ge2a$ and let $(\rho,\lambda)$ satisfy the
conditions of \cref{lem:tail} for $W_0=a^2/q_0^2$. For every odd $q\ge q_0$ prime to $a$ there is
$\lambda_q$ such that $(\rho,\lambda_q)$ satisfies them for $W=a^2/q^2$ and
\[
  \aconst_\sigma(a,q)-\bconst_q(\rho,\lambda_q)\ \ge\ \aconst_\sigma(a,q_0)-\bconst_{q_0}(\rho,\lambda)
  +\frac{10}{27}\log\frac q{q_0}.
\]
For even $a$, $\aconst_\sigma$ stands for the lower bound $\aconst_-$ (with $q^2+a^2$).
\end{proposition}

\begin{proof}
Take $\lambda_q=\lambda-\log(q/q_0)$; by \cref{prop:shift}, $\bconst_q(\rho,\lambda_q)\le
\bconst_{q_0}(\rho,\lambda)-\frac12\log(q/q_0)$. On the arithmetic side only the terms
$\frac13\log q-\Dk\log Y$ with $Y=q^2-\sigma a^2$ (or $q^2+a^2$) depend on $q$. The function
$s\mapsto\log(e^{2s}\mp a^2)$ has derivative $2/(1\mp a^2e^{-2s})\le2/(1-\frac14)=\frac83$ for
$e^s\ge2a$, so $\log Y(q)-\log Y(q_0)\le\frac83\log(q/q_0)$ and
\[
  \aconst_\sigma(a,q)-\aconst_\sigma(a,q_0)\ge\Bigl(\frac13-\frac83\Dk\Bigr)\log\frac q{q_0}.
\]
Since $\Dk\le\frac{25}{144}$, $\frac12+\frac13-\frac83\Dk\ge\frac56-\frac{25}{54}=\frac{10}{27}$.
\end{proof}

\begin{theorem}\label{thm:rows}
For each row $(a,\sigma,q_0,\kappa)$ of \cref{tab:constants} there are a comparison density $\rho$
and $\lambda$ satisfying the conditions of \cref{lem:tail} at $W_0=a^2/q_0^2$ with
$\aconst_\sigma(a,q_0)-\bconst(\rho,\lambda)>m_0$, where $m_0=\frac{13}{1000},\frac1{200},
\frac1{625},\frac3{2500},\frac7{10000},\frac7{10000}$ for the rows in the order of the table. For
$a=4,8$ the arithmetic constant is $\aconst_-$ (the one with $q^2+a^2$ and $c_2=C_e(\kappa)$).
\end{theorem}

\begin{proof}
This is the finite computation of Appendix~\ref{app:decay-certificates}, carried out with exact
rational arithmetic by the archived programs and with Arb balls by the second implementation.
\end{proof}

\begin{proof}[Proof of \cref{thm:main}]
Let $(a,\sigma)$ and $q\ge q_0(a,\sigma)$ be as in the statement, and $\kappa$ the ratio of the
corresponding row of \cref{tab:constants}; for $a=4,8$ take the row with $\sigma=-1$ for both signs,
which is legitimate because the bound of \cref{thm:decay} depends on $\abs w$ only and
$\aconst_+\ge\aconst_-$ for even $a$ (\cref{rem:worse-sign}). By \cref{thm:rows} and
\cref{prop:propagation}, $\bconst_q(\rho,\lambda_q)<\aconst_\sigma(a,q)$. By \cref{thm:arithmetic,thm:decay}
the first two hypotheses of \cref{prop:criterion} hold along $\mathcal N=\{n:T_n\ne0\}$. By
\cref{thm:nonvanishing,prop:local-dets}, for every $r\in\Q$ we have $T_p(r)\ne0$ for all large
primes $p$, since $q\ge14a$; in particular $T_n\ne0$ for infinitely many $n$. \Cref{prop:criterion}
gives $\xi\notin\Q$. For $-a/q$ use that $\chi_2$ and $\Ti_2$ are odd.
\end{proof}

\subsection{Status of the individual steps}\label{subsec:step-status}

We use three labels. \emph{Proved}: an elementary argument, given in full, of a kind we regard
as routine. \emph{Candidate}: an argument given here in full that has not yet been checked by
anyone independent of the AI systems involved. \emph{Certified (AI-internal)}: a finite
computation whose result was reproduced by a second implementation written without reference to
the first, both by AI systems; this is not an independent verification. \Cref{thm:main} as a
whole is a candidate.
\begin{center}
\small
\begin{tabular}{@{}p{0.44\textwidth}p{0.48\textwidth}@{}}
\toprule
Step & Status\\
\midrule
Moment sums, pole series, change of basis, criterion
(\cref{lem:moment-sums,lem:pole-series,lem:change-of-basis,prop:criterion}) & proved; the final
deduction is formalised in Lean, conditionally on the three estimates\\
Prime $2$, odd $a$ (\cref{thm:two-odd}) & candidate (new); entry bounds checked exactly on
$24\,960$ entries, finite forms consistent with exact valuations\\
Prime $2$, even $a$ (\cref{thm:two-even}) & candidate (new for $e\ge2$); consistent with exact
valuations at small $h$\\
Odd primes dividing $a$ (\cref{prop:p-divides-a}) & candidate (new); consistent with exact
valuations\\
Primes dividing $q(q^2-\sigma a^2)$ (\cref{sec:exceptional}) & candidate (adapted from Krisman and
the numerator-$2$ preprint)\\
Ordinary primes (\cref{thm:ordinary}) & candidate (adapted); finite bounds consistent with exact
valuations\\
Profile integrals $\Iprof_{7/6}$, $\Iprof_{5/4}$ & certified (AI-internal)\\
Size bound (\cref{thm:decay}) & candidate (adapted from Krisman)\\
Six strict margins (\cref{thm:rows}) & certified (AI-internal)\\
Propagation in $q$ (\cref{prop:propagation}) & proved (one-line calculus), conditional on the
candidate bounds it propagates; checked numerically\\
Nonvanishing mechanism (\cref{thm:nonvanishing}) & candidate (adapted); predicted valuation
confirmed at $p=13,17$\\
Local determinants (\cref{prop:local-dets}) & certified (AI-internal)\\
\bottomrule
\end{tabular}
\end{center}

\begin{remark}[points for scrutiny]\label{rem:scrutiny}
In our judgement the steps most in need of an independent check are: (i) the bookkeeping of
derivative columns and of the middle factor in the proof of \cref{thm:two-even}, in particular the
pole cost $e(2k+1)$ and the site cost $2em$ for $e\ge2$, on which the constants
$C_2(5/4)=-337/112$ and $C_3(5/4)=-9/2$ rest; (ii) the forward-difference analysis of
\cref{thm:two-odd}, in particular the cancellation of the term $\nu$ against the prime $2$ in
$\log(q^2-\sigma a^2)$; (iii) the rank-one argument of \cref{prop:p-divides-a}; (iv) the comparison
of the exact class data with the limiting profile in \cref{prop:profile}, on which the use of
midpoint sums as lower bounds depends; (v) the reduction of the local functionals in Step~5 of the
proof of \cref{thm:nonvanishing}; (vi) the implementations of the certificates in
Appendix~\ref{app:certificates}, which were written by AI systems. The margins (between $0.0007$ and
$0.0139$ per $h^2$) are small compared with the individual terms of \cref{tab:constants}, so an
error of a few thousandths in any constant would invalidate the corresponding row.
\end{remark}

\section{Limitations}\label{sec:limitations}

This section lists what the thresholds of \cref{thm:main} do and do not say; nothing in it is used in
the proof. The statements concern the specific estimates of this paper, not the determinants
themselves, nor the (ir)rationality of any number below a threshold.

\emph{The thresholds are not optimal.} They are the values at which the certified margin
$\aconst-\bconst$ of \cref{tab:constants} becomes positive with the comparison densities we
stored, for a mesh of $1/240$ and the two ratios $7/6$ and $5/4$. For $\chi_2(3/q)$ with
$\kappa=7/6$ the archived exact verifier finds the unfavourable gap $+0.0015$ per $h^2$ at $q=127$
with its stored density; we have not proved that no density does better there, nor examined other
pole ratios systematically. The rows $a=4$ and $a=8$ treat both signs together with the smaller
arithmetic constant; a size bound that sees the sign of $w$ \cite[Remark~6.7]{Krisman2026chi2}
could lower these thresholds for $\chi_2$.

\emph{Small margins.} The smallest margins ($7\cdot10^{-4}$ per $h^2$ for $\Ti_2(5/1609)$ and
for $a=8$) are of the order of $10^{-3}$ of the constants involved. The finite computations are
exact or enclosed in ball arithmetic, but the asymptotic estimates that feed them are candidates
(\cref{subsec:step-status}); a flaw of relative size $10^{-3}$ in any of them would remove a row.

\emph{Other numerators and the cases $a=1,2$.} The method applies to any numerator $a$ for
which the margin is positive; we did not search for further numerators. The numerators $1$ and $2$
are treated in~\cite{Krisman2026chi2,AxiomForge2026numeratorTwo}. Fractions $a/q$ with $q$ even are not
covered: the construction uses that $y=q/a$ and $w$ are $2$-adic units or have a known $2$-adic
valuation.

\emph{Comparison with the literature.} The range $q>54e^4a^2$ attributed to the criterion
of~\cite{DHK2022Acta} comes from a second-hand reading of its constants, not re-derived here
(\cref{subsec:background}); we make no claim of priority beyond the explicit thresholds of
\cref{thm:main}.


\appendix
\section{The finite computations}\label{app:certificates}

Four kinds of finite computations enter the proof. Each was first carried out by the code that
accompanies the candidate arguments, and then recomputed by a second implementation written
from the definitions in this paper (the scripts in the directory \texttt{checks/} of the
manuscript sources, see the section on reproducibility). Both implementations were written by AI
systems; agreement between them is what we call \emph{certified (AI-internal)}.

\subsection{The profile integrals}\label{app:profile-certificate}

For each $\kappa\in\{7/6,5/4\}$ and $\delta=1/2000$, the breakpoints $\mathfrak B\cap[\delta,2\kappa]$
of \cref{lem:convexity}, together with $\delta$ and $2\kappa$, cut $[\delta,2\kappa]$ into
$5999$ intervals. At the midpoint $x$ of each interval the class types, their capacities and $N(x)$
are rational, and $J(x)$, $R(x)$ are values of finite linear programmes with rational data. The
archived implementation fills the cheapest slots greedily after pre-filling the jet levels below
$\floor{1/x}-2$, which is legitimate because every slot of those levels is cheaper than every slot
at level $\floor{1/x}-1$ or above, and the levels below $\floor{1/x}-1$ have total capacity less
than one. The second implementation does not pre-fill: it finds by threshold search the least
integer cost $c^*$ such that the slots of cost at most $c^*$ have capacity at least one, using
closed forms, and fills up to $c^*$. All arithmetic is exact. Both give
\[
  \sum_{\text{intervals}}(b-a)\,B\bigl(\tfrac{a+b}2\bigr)-2\delta=
  \begin{cases}-2.11316413653\ldots&\kappa=7/6,\\
  -2.35563983202\ldots&\kappa=5/4,\end{cases}
\]
in agreement to all digits; the stated constants $\Iprof_{7/6}=-2113164137\cdot10^{-9}$ and
$\Iprof_{5/4}=-2355639833\cdot10^{-9}$ are rational numbers below the exact sums (the integral
is rounded down at the ninth decimal).

\subsection{The comparison densities and the six margins}\label{app:decay-certificates}

The six densities $\rho$ are step functions on the mesh $1/240$ with cell masses $m_i/24\,000\,000$,
$m_i\in\Z_{\ge0}$ and $\sum m_i=24\,000\,000$, cap $m_i\le200\,000$ (that is, $\rho\le2$), stored in \texttt{data/comparison.json} together with a rational $\lambda$
for each row. They were found by numerical optimisation; only their stored values matter.

The quantities in $\bconst(\rho,\lambda)$ are evaluated as follows.
\begin{itemize}
\item $E(\rho)$ is a finite sum over pairs of cells of
$\iint_{\text{cell}\times\text{cell}}\log\abs{x-y}=g(d+s)+g(d-s)-2g(d)$, with $s=1/240$, $d$ the
difference of the left endpoints, and $g(t)=\frac12t^2\log\abs t-\frac34t^2$.
\item $G$ is a finite sum of terms $c_j\psi(x-a_j)$ with rational $c_j$ and $a_j$ on the grid
$\frac1{24000}\Z$, plus $x\log W$. On each grid cell of length $\ell=1/24000$, $G$ is at most the
larger of its endpoint values plus $\sum_j\abs{c_j}\epsilon_j$, where $\epsilon_j$ bounds the
deviation of $\psi(x-a_j)$ from its chord: $\ell^2/(8\,\mathrm{dist})$ if $a_j$ is at distance
$\mathrm{dist}\ge\ell$ from the cell, and $\ell/e$ if $a_j$ is an endpoint (because
$\abs{t\log(t/\ell)}\le\ell/e$ on $[0,\ell]$).
\item The integral $2\int(G-\lambda)_+$ is bounded by $2\ell$ times the sum of the positive parts
of these cell bounds minus $\lambda$, and the tail conditions of \cref{lem:tail}
($1+\log W<0$, $\log W<\lambda$) are checked directly.
\end{itemize}
The archived implementation uses fixed-point logarithms with denominator $2^{100}$ and explicit
error budgets (and the cruder constant $\ell/2$); the second uses ball arithmetic (Arb, $200$ bits)
for every logarithm. They give, with decimals rounded in the safe direction (second implementation),
\begin{center}
\small
\begin{tabular}{@{}lllll@{}}
\toprule
row & $E(\rho)$ & $\bconst\le$ & $\aconst\ge$ & $\aconst-\bconst\ge$\\
\midrule
$3$, $131$, $\chi_2$ & $-2.124517762$ & $-3.187821144$ & $-3.173875358$ & $0.013945786$\\
$3$, $257$, $\Ti_2$ & $-2.154911545$ & $-3.546389753$ & $-3.541276253$ & $0.005113500$\\
$4$, $995$, both & $-2.177748279$ & $-4.083454335$ & $-4.081754611$ & $0.001699724$\\
$5$, $829$, $\chi_2$ & $-2.171320186$ & $-3.879559185$ & $-3.878258885$ & $0.001300300$\\
$5$, $1609$, $\Ti_2$ & $-2.180847665$ & $-4.212674322$ & $-4.211960841$ & $0.000713481$\\
$8$, $12061$, both & $-2.189990835$ & $-4.986095509$ & $-4.985336979$ & $0.000758530$\\
\bottomrule
\end{tabular}
\end{center}
which exceed the margins $m_0$ of \cref{thm:rows}; the archived verifier asserts the same margins in
exact rationals. The arithmetic constant uses the rational lower bounds $\Iprof_\kappa$ of
Appendix~\ref{app:profile-certificate}. For the two power-of-two rows the constant $c_2=C_e(\kappa)$
is the minimum of the quadratic $f_e$ of \cref{thm:two-even}, computed exactly ($-337/112$ at
$x=27/28$ for $e=2$ and $-9/2$ at $x=1$ for $e=3$, both at $\kappa=5/4$), and the logarithms of $2$ are
rounded in the direction of the sign of $c_2$. The second implementation also evaluated, for each row, the same density at
larger denominators ($q+2$, $q+12$, $2q+1$, $10q+1$, $1000q+1$, $\lambda$ shifted as in
\cref{prop:shift}) and found $B$ and the margin within the bounds of \cref{prop:shift,prop:propagation}.

\subsection{The local determinants}\label{app:local-certificates}

For both ratios the entries of $(Q-1)^4L_{s,e}(Q)$ are polynomials in $Q$ with rational
coefficients. The archived implementation computes the six determinants symbolically (SymPy). The
second evaluates $(Q-1)^{4t}\det L_{s,e}(Q)$ exactly at sufficiently many rational values of $Q$
(python-flint), interpolates, removes the factors $Q$ and $Q-1$, and compares with the stored
primitive integer polynomials: all six agree, with the exponents $(a_{s,e},b_{s,e})$ of
\cref{prop:local-dets}. The leading-term domination
$\sum_{i\ge1}\abs{\pi_{\deg-i}/\pi_{\deg}}R^{-i}<1$ at $R=83^2/4$, which also proves
nonvanishing for $\abs Q\ge R$, holds with left-hand sides between $4\cdot10^{-4}$ and $5.5\cdot10^{-3}$.

\subsection{Arithmetic entry bounds}\label{app:entry-certificates}

The archived program checks the anchored forward-difference entry bounds of
\cref{prop:two-odd-entries} exactly on $24\,960$ entries, for numerators $3,5,7,9$ and several
denominators, both signs, $h\le24$; and the exact quadratic minimisers $C_e(7/6)=-140/99$, $-227/84$,
$-614/153$ for $e=1,2,3$. The second implementation checks, entry by entry, the same bounds in
the finite-size comparison of Appendix~\ref{app:consistency}.

\section{Consistency checks at small sizes}\label{app:consistency}

The checks below compare the statements of the paper with exact computations of $T$ for small
$h$. They cannot confirm asymptotic statements, whose error terms are of the same size as the
differences observed, but they would have exposed errors in the finite forms of the bounds. The
polynomials $T$ were computed exactly by an implementation of the definitions in
\cref{sec:pencils} (python-flint rationals, determinants interpolated at $h+1$ rational points;
the script does not import the code that produced the candidate certificates). Each row of
\cref{tab:constants}, with both signs for $a=4,8$, was run for $n=1,2,\dots$ up to $h=32$--$56$; in
addition the numerators $9$ (with $q=35$, $\Ti_2$) and $12$ (with $q=1001$, $\chi_2$), which are
not rows of the theorem, were run to exercise $3\mid a$ with $9\mid a$ and $e=2$ with $a_0=3$.
No finite bound failed in any of the runs (\path{logs/experiments/runs/}).

\subsection{The pole series}
\Cref{lem:pole-series} was tested on five entries $\mu_\xi(\varphi u^s)$ with the series
$\sum_mw^m(\varphi u^s)'(m)$ summed directly (exact rational derivatives, $\xi$ summed with $700$
digits), for the six threshold rows and both signs of $a=4,8$ at $n=1$. The largest relative
discrepancy is $2\cdot10^{-664}$.

\subsection{The prime $2$}
For odd $a$ every entry bound of \cref{prop:two-odd-entries} was checked in the anchored bases
(for example $1936$ entries for $a=5$, $q=829$, $h=44$), with no failure, and the exact
$\vG2(T)$ lies above the finite bound. For $a=4$, $q=995$ and $h=56$, $\vG2(T)=-9156$ against the
finite bound $-9168$; for $a=8$, $q=12061$ the finite bound is attained exactly for every
$h=4n\le48$ ($-63$ at $h=4$, $-10184$ at $h=48$). Since the leading terms are
$C_e(\kappa)h^2-O(h\log h)$ with a logarithmic error of the same size as the differences observed,
we do not tabulate ratios.

\subsection{Odd primes}
For each run and each $n$ the primes dividing $a$, $q$ and $q^2-\sigma a^2$ were tested against the
finite form of the relevant bound: \cref{prop:p-divides-a} (for example for $a=3$, $q=131$, $h=48$,
$\vG3(T)/h^2=-1.2812$ against $-(2\kappa-1)=-1.3333$), \cref{prop:p-divides-q} (the ratio
$\vG p(T)/(fh^2)$ is $0.3332$ for $p=199\mid995$ at $h=56$, and below $1/3$ by amounts
compatible with $O(h\log h)$ for $p=5\mid995$: $0.3084$) and \cref{prop:p-divides-Y}. No bound
failed. The ordinary primes ($p>2K$ and $p^2>2K$ via the jet and pole bounds) were checked
in the earlier numerator-$2$ audit; for the numerators of this paper only the exceptional primes
were tested at each size.

\subsection{Nonvanishing}
\Cref{thm:nonvanishing} predicts $v_p(\Delta(0))=-pt(\gamma+2-t)+t$ at $h=tp$. Exact computation of
$\Delta(0)$ gives $v_{13}=-152$ for $(a,q,\sigma)=(4,995,1)$, $(3,257,-1)$ and $(8,12061,-1)$ with
$h=52$ and $v_{17}=-200$ for $(5,829,1)$ with $h=68$, and $v_{17}=-300$ for $(3,131,1)$ with
$\kappa=7/6$ and $h=102$, in each case as predicted. These are checks of the whole mechanism of
\cref{sec:nonvanishing}, including the reduction of the local functionals, at one prime per row.

\subsection{Size}
The size bound \cref{thm:decay} was not compared with exact determinants for these rows: the
margins are tiny, and the finite sizes available are too small for $O(\log h/h)$ corrections to be
negligible. The comparison of the six exponents with those obtained in the numerator-$2$ work, of
the same order, is the only sanity check.


% back matter starts on a new page
\clearpage
\afbacksection{Reproducibility and formal verification}
\sloppy

\emph{Status of the result.} \Cref{thm:main} is a candidate. No human expert independent of the
AI systems that produced it has reviewed the argument. Each step carries one of the labels of
\cref{subsec:step-status}: \emph{proved} (routine elementary arguments), \emph{candidate}
(argued in this paper, not yet independently verified) or \emph{certified (AI-internal)} (a finite
computation reproduced by a second implementation, both written by AI systems). None of this is
an independent verification.

\emph{Computations.} Paths are relative to the directory of this contribution (AF-0004) in
the AxiomForge archive (\url{https://github.com/axiomforge-math/mathematics}); every command, its
output and its outcome are recorded in \path{logs/verification.md}. The scripts used for
Appendices~\ref{app:certificates} and~\ref{app:consistency} (the directory \texttt{checks/}
mentioned there) are in \path{logs/experiments/scripts/}, and the step densities and integer
polynomials (the directory \texttt{data/}) are the following files of
\path{logs/experiments/certificates/}:
\begin{itemize}
\item \path{comparison_densities.json} (the six densities and $\lambda$),
\item \path{local_polynomials.json} (the six integer polynomials).
\end{itemize}
The scripts need Python~3.11 or later with \texttt{python-flint} and \texttt{mpmath}; none of them
imports the code that produced the candidate certificates. They are:
\begin{itemize}
\item \path{pencil.py}: exact implementation of $\mu_X$, $\Delta$, $T$ (used by the three checks
below);
\item \path{profile.py}: the profile integrals of \cref{thm:ordinary} (exact rational arithmetic);
\item \path{decay_ball.py}: the rigorous bounds for $\bconst$ and $\aconst$ and the six margins of
Appendix~\ref{app:decay-certificates} (Arb ball arithmetic), with an option for larger $q$;
\item \path{local_poly.py}: the local determinants of \cref{prop:local-dets};
\item \path{series_identity.py}, \path{finite_bounds.py} and \path{global_prime_index.py}: the
consistency checks of Appendix~\ref{app:consistency}.
\end{itemize}
The archived programs that first produced the certificates (exact rational verifiers for the
profile integrals, the comparison densities and the entry bounds, and a symbolic computation of
the local determinants) were also re-executed and reproduce their logs; they are not part of this
package.

\emph{Formal verification.} The Lean~4 package in \path{lean/} (Lake package
\texttt{axiomforge-af0004}, root module \texttt{AxiomForge.AF0004}, namespace
\texttt{AxiomForge.LegendreChiGeneral}) is adapted from the package of the numerator-$2$
preprint. It contains: definitions of $\chi_2$ and $\Ti_2$ as series, with their oddness; an
explicit definition of the polynomials $T_n$ of \cref{sec:pencils}; the irrationality
criterion (\cref{prop:criterion}) in several forms, proved unconditionally; the exact evaluation of
the constants $c_2$ of \cref{thm:three-estimates}; and the six strict margins of \cref{thm:rows}
for the rational bounds of \cref{tab:constants}, proved by \texttt{norm\_num} (no
\texttt{native\_decide}). The main theorems are \emph{not} formalised. The package proves only the
final deduction, conditionally: the theorems whose names end in \texttt{\_of\_estimates} derive the
irrationality of $\chi_2(\pm a/q)$ and $\Ti_2(\pm a/q)$ for the six rows from a hypothesis that
packages the arithmetic estimate, the size estimate and the nonvanishing at prime indices for these
polynomials, with rounded rates. These hypotheses are exactly the unproved content of
\cref{sec:prime-two,sec:exceptional,sec:ordinary,sec:decay,sec:nonvanishing}; the Lean theorems are
therefore not formalisations of \cref{thm:main}. The formalization of the three estimates is in
preparation. The package builds with \texttt{lake exe cache get \&\& lake build};
\texttt{\#print axioms} reports only \texttt{propext}, \texttt{Classical.choice} and
\texttt{Quot.sound} for each audited declaration (\path{lean/AXIOMS.txt}), and the sources contain
no \texttt{sorry}, \texttt{admit}, \texttt{axiom} declaration or \texttt{native\_decide}. No
independent human review of this work has taken place.

\afbacksection{Acknowledgements}

The construction studied here starts from the work of Shawn Krisman~\cite{Krisman2026chi2,
Krisman2026chi2small}, whose determinants, criterion, reflection argument, size bound and
nonvanishing strategy are reused and adapted throughout, and continues the AxiomForge preprint on the
numerator $2$~\cite{AxiomForge2026numeratorTwo}; we have tried to indicate each such use at the place
where it occurs.

\afbacksection{AI-assisted research disclosure}

\emph{Exploration and candidate proofs.} According to the owners of AxiomForge, the original
exploration, the choice of the two constructions, the candidate proofs and the first versions
of the certificate programs were produced by GPT-6.2 XHigh (OpenAI), used through its web
interface.

\emph{Reconstruction, audit and manuscript.} Claude Sonnet 5.5, Claude Fable 5.2 and Claude Opus
5.5 (Anthropic), working under the direction of AxiomForge, reconstructed the arguments from the
candidate material, rewrote the proofs, wrote the second implementations of the finite computations
(Appendix~\ref{app:certificates}), carried out the small-size checks (Appendix~\ref{app:consistency}),
wrote this manuscript, and wrote the Lean package described in the section on reproducibility.
Their reading of the arithmetic lemmas is an AI-internal review: it was made by the same family of
systems that wrote the text, and is not an independent verification.

\emph{Human role and review.} The owners of AxiomForge chose the problem, directed the work and
decided on its publication as a preprint. No human expert independent of the AI systems has
reviewed the mathematics. The checks reported here were performed by AI systems and are not an
independent certification; AI involvement in a step should not be read as evidence of its
correctness.

\emph{Responsibility.} AI systems are not authors of this work and cannot take responsibility
for it. AxiomForge takes full responsibility for the entire content of this manuscript,
including every part produced with AI assistance, and for any errors it may contain. Reports
of errors are welcome at \href{mailto:axiomforge.math@gmail.com}{axiomforge.math@gmail.com};
confirmed errors are corrected openly in a new version, recorded in the changelog of this
contribution, with credit to whoever reported them.


\clearpage
\bibliographystyle{amsplain}
\bibliography{references}

\end{document}
